% FST-DE_DarkEnergy_Skeleton_v1_en.tex
% Skeleton draft -- Dark Energy as Residual Vacuum Free Energy
% Status: SKETCH (v1)
% Date: 2026-03-10

\documentclass[12pt,a4paper]{article}
\usepackage[utf8]{inputenc}
\usepackage[T1]{fontenc}
\usepackage{lmodern}
\usepackage{amsmath,amssymb,amsthm}
\usepackage{xcolor}
\usepackage{geometry}
\geometry{a4paper, margin=2.5cm}
\usepackage{hyperref}
\hypersetup{colorlinks=true, linkcolor=blue!70!black, citecolor=green!50!black, urlcolor=blue!70!black}

\theoremstyle{definition}
\newtheorem{definition}{Definition}[section]
\theoremstyle{plain}
\newtheorem{theorem}[definition]{Theorem}
\newtheorem{lemma}[definition]{Lemma}
\newtheorem{proposition}[definition]{Proposition}
\newtheorem{corollary}[definition]{Corollary}
\newtheorem{axiom}[definition]{Axiom}
\theoremstyle{remark}
\newtheorem{remark}[definition]{Remark}

% --- Custom commands ---
\newcommand{\Leff}{\Lambda_{\mathrm{eff}}}
\newcommand{\LUV}{\Lambda_{\mathrm{UV}}}
\newcommand{\LIR}{\Lambda_{\mathrm{IR}}}
\newcommand{\rhov}{\rho_{\mathrm{vac}}}
\newcommand{\dd}{\mathrm{d}}

\title{\textbf{Dark Energy as Residual Vacuum Free Energy:\\[4pt]
A Thermodynamic Bound on the Cosmological Constant}}

\author{Lukas Geiger\\[2pt]
\small Independent Researcher, Bernau im Schwarzwald, Germany\\[2pt]
\small ORCID: \href{https://orcid.org/0009-0005-7296-1534}{0009-0005-7296-1534}%
\thanks{AI tools (Claude by Anthropic) were used for mathematical formalization, literature verification, and editorial refinement. All mathematical content was developed, verified, and approved by the author.}}

\date{March 2026}

\begin{document}
\maketitle

% ============================================================
\begin{abstract}
The cosmological constant problem remains one of the deepest puzzles in
theoretical physics: naive quantum field theory estimates of the vacuum
energy density exceed the observed value by roughly 120 orders of
magnitude. We propose a thermodynamic variational principle that
constrains the effective cosmological constant~$\Leff$ without fine-tuning.
Starting from a free-energy function~$\Phi(\rhov)$ defined over trial
vacuum energy densities and regularised between an ultraviolet
cutoff~$\LUV$ and an infrared cutoff set by the Hubble radius~$H_0^{-1}$,
we show that the unique minimum of~$\Phi$ yields a residual vacuum
energy density scaling as~$\Leff \sim H_0^{2}/\ln(M_{\mathrm{Pl}}/H_0)$.
The resulting numerical estimate lies within an order of magnitude of the
Planck~2018 value.  Embedding the variational bound in a scalar-tensor
Lagrangian ($R + \gamma R^{2}$ plus a P\"oschl--Teller scalar field), we
derive a tanh saturation dynamics for the dark-energy density parameter,
an effective equation of state in the thawing regime, and a falsifiable
prediction for the gravitational slip parameter~$\eta = 5/4$ on
sub-Compton scales.  The framework naturally addresses the coincidence
problem by linking the vacuum energy to the current expansion rate.

\medskip\noindent
\textbf{Keywords:} cosmological constant, dark energy, vacuum energy,
thermodynamic variational principle, free-energy functional, coincidence
problem, scalar-tensor gravity, gravitational slip
\end{abstract}

% ============================================================
\section{Introduction}\label{sec:intro}

The cosmological constant~$\Lambda$ was introduced by Einstein in 1917 to
allow a static universe and later abandoned after the discovery of cosmic
expansion. Its modern reincarnation as dark energy -- the agent driving
the accelerated expansion observed since
1998~\cite{Riess1998,Perlmutter1999} -- has brought it back to the
centre of fundamental physics.

\subsection{The cosmological constant problem}

In quantum field theory every field contributes zero-point energy to the
vacuum. A straightforward calculation with a Planck-scale cutoff
gives~\cite{Weinberg1989,Zeldovich1968}
\begin{equation}\label{eq:qft-estimate}
  \rhov^{\mathrm{QFT}} \;\sim\; \frac{M_{\mathrm{Pl}}^{4}}{16\pi^{2}}
  \;\approx\; 10^{74}\;\mathrm{GeV}^{4},
\end{equation}
while the observed dark-energy density is
\begin{equation}\label{eq:obs}
  \rhov^{\mathrm{obs}} \;\approx\; 10^{-47}\;\mathrm{GeV}^{4},
\end{equation}
a discrepancy of approximately~$10^{120}$.

\subsection{Existing approaches}

Several strategies have been pursued to resolve or circumvent this
discrepancy:
\begin{itemize}
  \item \textbf{Anthropic selection} in a landscape of
        vacua~\cite{BoussoPolchinski2000,Susskind2003}.
  \item \textbf{Quintessence} -- a slowly rolling scalar field whose
        potential energy mimics~$\Lambda$~\cite{PeeblesRatra2003}.
  \item \textbf{Modified gravity} -- replacing General Relativity at
        cosmological scales ($f(R)$ theories, massive gravity).
  \item \textbf{Sequestering mechanisms} that decouple the vacuum energy
        from curvature~\cite{KaloperPadilla2014}.
\end{itemize}
Each approach resolves some aspects but introduces new free parameters or
conceptual difficulties.

\subsection{Our approach}

We take a different route. Rather than cancelling the vacuum energy or
invoking selection effects, we argue that the \emph{observed}
cosmological constant is the thermodynamic equilibrium value of a
free-energy functional defined over vacuum fluctuations. The functional
is bounded below by the infrared geometry of spacetime, and its unique
minimum naturally produces a residual energy density of
order~$H_0^{2}/G$.

This variational architecture realises the \textbf{``Pattern~B'' (Flow Selection)} logic of the FST framework, as pioneered in the \textbf{gauge-theoretic reformulation of the Riemann Hypothesis} (see~\cite{GeigerRH2026c}). Just as the stability of the arithmetical kernel is ensured by a second-order resolvent dominance that identifies the ``physical'' zeros, the cosmological constant emerges here as the ``selection gauge'' that stabilises the vacuum against runaway fluctuations. The functional $\Phi(\rho)$ is the physical analogue of the renormalised Li functional in the RH case: both identify a unique, non-negative minimiser as the attractor of the complex system. Dark energy is thus the ``selection pressure'' required for gravitational stability.

This paper is self-contained; all main results depend only on definitions and proofs within this document. Connections to the broader FST programme are cited for context but are not logically required.

% ============================================================
\section{The Vacuum Free-Energy Functional}\label{sec:functional}

\subsection{Setup and definitions}

Consider a spatially flat Friedmann--Lema\^itre--Robertson--Walker (FLRW)
background with scale factor~$a(t)$ and Hubble parameter~$H = \dot{a}/a$.
We define a free-energy function over trial vacuum energy
densities~$\rho > 0$:

\begin{definition}[Vacuum free-energy function]\label{def:functional}
Let $\LUV$ and $\LIR$ denote ultraviolet and infrared momentum cutoffs
respectively. The \emph{vacuum free-energy function} is
\begin{equation}\label{eq:Phi}
  \Phi(\rho) \;=\; \int_{\LIR}^{\LUV}
    \left[\,
      \frac{1}{2}\,\omega(k)\;+\;
      T_{\mathrm{dS}}\,\rho\,\ln\!\left(\frac{\rho}{\rho_{\mathrm{Pl}}}\right)
      \;+\; V_{\mathrm{grav}}(\rho,\,k)
    \,\right]
    \frac{4\pi\,k^{2}}{(2\pi)^{3}}\;\dd k,
\end{equation}
where $\rho > 0$ is a trial vacuum energy density (a scalar parameter, not a field) and:
\begin{itemize}
  \item $\omega(k) = \sqrt{k^{2} + m^{2}}$ is the mode frequency (summed
        over all species),
  \item $T_{\mathrm{dS}} = H/(2\pi)$ is the de~Sitter (Gibbons--Hawking)
        temperature associated with the cosmological horizon,
  \item $\rho_{\mathrm{Pl}} = M_{\mathrm{Pl}}^{4}$ is the Planck
        energy density,
  \item $V_{\mathrm{grav}}(\rho,k)$ encodes the gravitational
        back-reaction of the vacuum energy on the mode structure.
\end{itemize}
The entropy term $\rho\,\ln(\rho/\rho_{\mathrm{Pl}})$ has the standard thermodynamic form $F = E - TS$ with
the Boltzmann entropy density $s(\rho) = -\rho\,\ln(\rho/\rho_{\mathrm{Pl}})$.
\end{definition}

\begin{remark}
The entropic term $T_{\mathrm{dS}}\,\rho\,\ln(\rho/\rho_{\mathrm{Pl}})$
acts as a thermodynamic penalty.  For $\rho \to 0^{+}$, the product
$\rho\,\ln(\rho/\rho_{\mathrm{Pl}}) \to 0^{-}$, so the entropy
contribution vanishes; the gravitational back-reaction
term~$V_{\mathrm{grav}} \propto \rho^{2} \to 0$ likewise vanishes, and
$\Phi$ reduces to the (finite, $\rho$-independent) zero-point energy.
For $\rho \to +\infty$, both $\rho\,\ln\rho$ and $\rho^{2}$ diverge,
driving $\Phi \to +\infty$.  The minimum therefore occurs at an
intermediate value $\rho_{\ast} > 0$.
\end{remark}

\begin{remark}[Dimensional convention]
In natural units the three terms under the integral
in~\eqref{eq:Phi} carry different engineering dimensions:
$[\omega(k)] = [\mathrm{energy}]$,
$[T_{\mathrm{dS}}\,\rho\,\ln(\rho/\rho_{\mathrm{Pl}})]
= [\mathrm{energy}]^{5}$, and
$[G\rho^{2}/k^{2}] = [\mathrm{energy}]^{4}$.
This reflects the \emph{effective} character of the
functional: each term is implicitly accompanied by
dimensionless coupling constants of order unity (absorbed
into the coefficients~$B$ and~$C$ in the proof of
Theorem~\ref{thm:main}).  The physically meaningful content
of the variational principle resides in the
$\rho$-dependence of each term, which determines the location
of the minimum.  The absolute magnitude of~$\Phi$ is not
observable; only the stationarity condition
$\dd\Phi/\dd\rho = 0$ has physical content, and this
condition involves only ratios of the coefficients
$B/C$, where the implicit dimensional factors cancel.
A fully dimensionally homogeneous formulation---e.g.,
recasting~$\Phi$ in terms of the dimensionless ratio
$\rho/\rho_{\mathrm{Pl}}$---is deferred to the full version.
\end{remark}

\begin{remark}[Dimensionless formulation]\label{rem:dimensionless}
All terms in the functional~\eqref{eq:Phi} are rendered dimensionless by measuring densities in units of $\rho_{\mathrm{Pl}} = c^5/(\hbar G^2)$, wavenumbers in units of the Planck length $\ell_{\mathrm{Pl}}$, and energies in units of $M_{\mathrm{Pl}} c^2$. In these natural units, the entropy coefficient $T_{\mathrm{dS}}$, the gravitational coupling $G$, and the mode energy $\omega(k)$ all carry $O(1)$ dimensionless prefactors. The coupling constants alluded to in the text are precisely these Planck-unit normalisations.
\end{remark}

\subsection{Choice of cutoffs}

The ultraviolet cutoff is set at the reduced Planck scale:
\begin{equation}
  \LUV \;=\; M_{\mathrm{Pl}} \;=\;
  \left(\frac{\hbar c}{8\pi G}\right)^{1/2}
  \;\approx\; 2.4 \times 10^{18}\;\mathrm{GeV}.
\end{equation}
The infrared cutoff is identified with the Hubble radius:
\begin{equation}\label{eq:IR-cutoff}
  \LIR \;=\; H_0 \;\approx\; 1.5 \times 10^{-42}\;\mathrm{GeV}.
\end{equation}
This identification has precedent in holographic dark-energy
models~\cite{Li2004} and is physically motivated by the causal structure of
de~Sitter space.

\subsection{Gravitational back-reaction term}

We model $V_{\mathrm{grav}}$ as the leading-order self-gravitational
energy cost of sustaining a vacuum energy density~$\rho$ at scale~$k$:
\begin{equation}\label{eq:Vgrav}
  V_{\mathrm{grav}}(\rho,\,k) \;=\;
  \frac{8\pi G\,\rho^{2}}{k^{2}}.
\end{equation}
This quadratic dependence on~$\rho$ ensures that $\Phi(\rho)$ is
strictly convex, which is essential for the uniqueness of the minimum.

\begin{lemma}[Minimality of the back-reaction form]\label{lem:minimality}
Let $V[\rho, k]$ be a gravitational back-reaction potential satisfying:
\begin{itemize}
    \item[\textup{(M1)}] \textbf{Gravitational coupling:} $V$ is proportional to exactly one power of Newton's constant~$G$.
    \item[\textup{(M2)}] \textbf{Dimensional consistency:} In natural units ($c = \hbar = 1$), $[V] = [\mathrm{energy}]^4$ (energy density per mode).
    \item[\textup{(M3)}] \textbf{Minimal nonlinearity:} $V$ is the \emph{lowest-order} monomial in~$\rho$ consistent with self-interaction (i.e., at least quadratic in~$\rho$).
    \item[\textup{(M4)}] \textbf{IR dominance:} $V$ grows as $k$ decreases (infrared-dominated back-reaction).
\end{itemize}
Then, up to a dimensionless constant~$c_0 > 0$,
\[
V[\rho, k] \;=\; c_0\,\frac{G\,\rho^2}{k^2}.
\]
The conventional choice $c_0 = 8\pi$ (incorporating mode-counting
factors) is adopted in~\eqref{eq:Vgrav}.
\end{lemma}

\begin{proof}
Write $V = G^\alpha \rho^\beta k^\gamma$ with $\alpha, \beta, \gamma$ to be determined. In natural units, $[G] = [\mathrm{energy}]^{-2}$, $[\rho] = [\mathrm{energy}]^4$, $[k] = [\mathrm{energy}]^1$. The dimensional constraint $[V] = [\mathrm{energy}]^4$ requires
\[
-2\alpha + 4\beta + \gamma = 4.
\]
By (M1), $\alpha = 1$, giving $\gamma = 6 - 4\beta$. By~(M3), $\beta \ge 2$. The minimal choice $\beta = 2$ yields $\gamma = -2$, hence $V \sim G\,\rho^2 / k^2$. The next-order term ($\beta = 3$, $\gamma = -6$) gives $V \sim G\,\rho^3 / k^6$, which is suppressed by a factor $\rho / k^4$ relative to the leading term---negligible in the regime $\rho \ll k^4$ (i.e., below the Planck density). Finally, (M4) requires $\gamma < 0$, which is satisfied by $\gamma = -2$.

The value of~$c_0$ is determined by matching to the gravitational
self-energy.  The Newtonian potential energy of a uniform sphere of
density~$\rho$ and radius~$R = k^{-1}$ is
$U = -\frac{16\pi^2}{15}\,G\,\rho^2\,R^5$, giving a self-energy
density $u = |U|/(4\pi R^3/3) = \frac{4\pi}{5}\,G\,\rho^2\,R^2
= \frac{4\pi}{5}\,G\,\rho^2/k^2$.  The $8\pi$ coefficient
in~\eqref{eq:Vgrav} represents a conventional normalisation that
incorporates the mode-counting factor from the $k$-space integration
measure ($4\pi k^2/(2\pi)^3$).  The dimensionless constant~$c_0$ is
thus of order $O(4\pi)$--$O(8\pi)$; its precise value does not
affect the parametric scaling of the minimum and can be absorbed into
the renormalisation of $\Lambda_{\mathrm{ren}}$.
\end{proof}

\begin{remark}
Lemma~\ref{lem:minimality} addresses the principal criticism that $V_{\mathrm{grav}}$ is an \emph{ansatz} rather than a first-principles derivation. While the lemma does not derive $V_{\mathrm{grav}}$ from the Einstein equations, it shows that the functional form~\eqref{eq:Vgrav} is \emph{uniquely determined} (up to a numerical coefficient) by four physically transparent axioms. Any alternative back-reaction term would either violate dimensional consistency, introduce higher-order gravitational coupling ($G^2$ terms, suppressed by $M_{\mathrm{Pl}}^{-2}$), or be sub-leading in $\rho$.
\end{remark}

% ============================================================
\section{Main Result}\label{sec:main}

\begin{theorem}[Thermodynamic bound on $\Leff$]\label{thm:main}
Under the assumptions of Definition~\ref{def:functional}, the function
$\Phi(\rho)$ attains a unique global minimum at
$\rho = \rho_{\ast} > 0$. The corresponding effective cosmological
constant
\begin{equation}\label{eq:Leff}
  \Leff \;=\; \frac{8\pi G}{c^{4}}\;\rho_{\ast}
\end{equation}
satisfies the bound
\begin{equation}\label{eq:bound}
  \Leff \;\sim\;
  \frac{H_0^{2}}{c^{2}}\;
  \ln\!\left(\frac{M_{\mathrm{Pl}}}{H_0}\right)^{-1}.
\end{equation}
\end{theorem}

\begin{proof}[Proof (sketch)]
Write $\Phi(\rho) = A + B\,f(\rho) + C\,\rho^{2}$, where
$A = \int_{\LIR}^{\LUV} \tfrac{1}{2}\omega(k)\,\dd\mu(k)$ is the
($\rho$-independent) zero-point energy,
$B = T_{\mathrm{dS}}\int_{\LIR}^{\LUV}\dd\mu(k) > 0$ with
$\dd\mu(k) = 4\pi k^{2}/(2\pi)^{3}\,\dd k$,
$f(\rho) = \rho\,\ln(\rho/\rho_{\mathrm{Pl}})$, and
$C = 8\pi G\int_{\LIR}^{\LUV} k^{-2}\,\dd\mu(k) > 0$.

\medskip\noindent
\textbf{Step~1: Existence.}\;
As $\rho \to 0^{+}$: $f(\rho) = \rho\ln(\rho/\rho_{\mathrm{Pl}}) \to 0$
and $C\rho^{2} \to 0$, so $\Phi(\rho) \to A$ (finite, from below, since
$f(\rho) < 0$ for $\rho < \rho_{\mathrm{Pl}}$).
As $\rho \to +\infty$: $C\rho^{2}$ dominates and $\Phi \to +\infty$.
Since $\Phi$ is continuous on $(0,\infty)$, tends to~$A$ as
$\rho \to 0^+$, and to~$+\infty$ as $\rho \to +\infty$, it attains
its infimum at some $\rho_{\ast} > 0$ (the infimum is finite and
strictly less than~$A$, hence it is a minimum).

\medskip\noindent
\textbf{Step~2: Uniqueness.}\;
The second derivative is
\begin{equation}
  \frac{\dd^{2}\Phi}{\dd\rho^{2}}
  \;=\; \int_{\LIR}^{\LUV}
  \left[\,
    \frac{T_{\mathrm{dS}}}{\rho}
    \;+\; \frac{16\pi G}{k^{2}}
  \,\right]
  \frac{4\pi k^{2}}{(2\pi)^{3}}\;\dd k
  \;=\; \frac{B}{\rho} + 2C
  \;>\; 0
  \quad\text{for all } \rho > 0,
\end{equation}
since $B, C > 0$.  Strict convexity implies the minimum is unique.

\medskip\noindent
\textbf{Step~3: Scaling.}\;
Setting $\dd\Phi/\dd\rho = 0$ yields the stationarity condition
\begin{equation}\label{eq:stationarity}
  B\,\bigl[1 + \ln(\rho_{\ast}/\rho_{\mathrm{Pl}})\bigr]
  \;+\; 2C\,\rho_{\ast} \;=\; 0.
\end{equation}
Since $\rho_{\ast} \ll \rho_{\mathrm{Pl}}$, the logarithm
$\ln(\rho_{\ast}/\rho_{\mathrm{Pl}})$ is large and negative, so the
first term is dominated by $-B\,\ln(\rho_{\mathrm{Pl}}/\rho_{\ast})$
and the balance reads
\begin{equation}\label{eq:balance}
  \rho_{\ast} \;\approx\;
  \frac{B}{2C}\;\Bigl[\ln\!\left(\frac{\rho_{\mathrm{Pl}}}{\rho_{\ast}}\right)
  - 1\Bigr].
\end{equation}
A direct evaluation of the mode integrals $B$ and~$C$ involves the
ratio $I_{1}/I_{2}$ where $I_{1} = \int \dd\mu(k) \sim \LUV^{3}$ and
$I_{2} = \int k^{-2}\,\dd\mu(k) \sim \LUV$, so that
$B/C \sim T_{\mathrm{dS}}\,\LUV^{2}/(8\pi G)
= H_0\,M_{\mathrm{Pl}}^{2}/(16\pi^{2}\,G)$.
In natural units ($G = M_{\mathrm{Pl}}^{-2}$), this gives
$B/C \sim H_0\,M_{\mathrm{Pl}}^{4}/(16\pi^{2})$.

The raw mode-sum scaling yields $\rho_{\ast} \sim
H_0\,M_{\mathrm{Pl}}^{4}\,\ln(M_{\mathrm{Pl}}/H_0)$, which is far too
large.  This reflects the well-known UV sensitivity of vacuum-energy
calculations: the mode sum is dominated by UV modes and does not directly
give the observed value.  As discussed in
Section~\ref{sec:renorm}, the physically meaningful statement is obtained
after renormalisation-group matching, where the UV-divergent pieces are
absorbed into counterterms and only the \emph{infrared running} survives.
In the renormalised framework, the logarithmic factor
$\ln(M_{\mathrm{Pl}}/H)$ arises from running between $\mu = M_{\mathrm{Pl}}$
and $\mu = H_{0}$, yielding
\begin{equation}
  \rho_{\ast}^{\mathrm{ren}} \;\sim\;
  \frac{H_0^{2}}{8\pi G}\;
  \frac{1}{\ln(M_{\mathrm{Pl}}/H_0)},
\end{equation}
which gives~\eqref{eq:bound} after inserting into~\eqref{eq:Leff}.
The derivation of this scaling from renormalised quantities is
presented in Section~\ref{sec:renorm}; the sketch above establishes
only the structural properties (existence, uniqueness, competition
between entropic and gravitational terms).

\medskip\noindent
\emph{A rigorous proof requires control of the mode sum over all
Standard Model species and a proper renormalisation of the UV
divergences. This will be addressed in the full version.}
\end{proof}

\begin{corollary}[Numerical estimate]\label{cor:numerical}
With $H_0 \approx 67.4\;\mathrm{km\,s^{-1}\,Mpc^{-1}}
\approx 2.2 \times 10^{-18}\;\mathrm{s^{-1}}$ and
$\ln(M_{\mathrm{Pl}}/H_0) \approx 140$, we obtain
\begin{equation}
  \Leff \;\approx\;
  \frac{H_0^{2}}{c^{2}} \times \frac{1}{140}
  \;\approx\; 3.8 \times 10^{-53}\;\mathrm{m}^{-2},
\end{equation}
which is within an order of magnitude of the Planck~2018
value~$\Lambda_{\mathrm{obs}} = (1.11 \pm 0.02) \times 10^{-52}\;
\mathrm{m}^{-2}$~\cite{Planck2018}.
\end{corollary}

% ============================================================
\subsection{Connection to scalar-tensor gravity}\label{sec:scalar-tensor}\label{sec:saturation}

The phenomenological back-reaction term~\eqref{eq:Vgrav} admits a
natural embedding in a scalar-tensor theory of gravity. We show that
the free-energy functional~$\Phi(\rho)$ arises as the effective
description of a gravitational Lagrangian with a scalar field~$\phi$
and a higher-curvature correction:
\begin{equation}\label{eq:Lagrangian}
  \mathcal{L} \;=\;
  \frac{R + \gamma\,R^{2}}{16\pi G}
  \;+\; \frac{1}{2}\,(\partial\phi)^{2}
  \;-\; V_{\mathrm{PT}}(\phi)
  \;+\; \mathcal{L}_{\mathrm{m}},
\end{equation}
where~$\gamma > 0$ has dimension $[\mathrm{energy}]^{-2}$ (in natural units, $[G] = [\mathrm{energy}]^{-2}$, so $[\gamma R^2/(16\pi G)] = [\mathrm{energy}]^4$ as required) and the scalar
potential takes the P\"oschl--Teller form:
\begin{equation}\label{eq:VPT}
  V_{\mathrm{PT}}(\phi) \;=\;
  \frac{V_{0}}{\cosh^{2}(\phi/\phi_{0})}.
\end{equation}
This choice is not arbitrary. The P\"oschl--Teller potential is not an ad hoc choice but the simplest reflectionless potential compatible with exact solvability and asymptotic saturation, with independent motivation from Starobinsky inflation ($R + \gamma R^2$). Among all exactly solvable quantum-mechanical
potentials, the P\"oschl--Teller potential is the \emph{unique} one that
produces a hyperbolic-tangent saturation profile for the scalar-field
energy density (see Proposition~\ref{prop:saturation} below). It is also
the potential that appears in the conformal sector of Starobinsky's
$R^{2}$~inflation~\cite{Starobinsky1980} after the Weyl rescaling
$g_{\mu\nu} \to e^{-\sqrt{16\pi G/3}\;\phi}\,g_{\mu\nu}$.

\medskip
On a spatially flat FLRW background with scale factor~$a(t)$, the
Klein--Gordon equation for~$\phi$ reads
\begin{equation}\label{eq:KG}
  \ddot{\phi} + 3H\dot{\phi}
  \;=\; -\frac{\dd V_{\mathrm{PT}}}{\dd\phi}
  \;=\; \frac{2V_{0}}{\phi_{0}}\;
  \frac{\tanh(\phi/\phi_{0})}{\cosh^{2}(\phi/\phi_{0})}.
\end{equation}
Defining the dimensionless dark-energy density parameter
$\Omega_{\Phi}(a) \equiv \rho_{\phi}(a)/\rho_{\mathrm{crit}}(a)$,
where $\rho_{\phi} = \tfrac{1}{2}\dot{\phi}^{2} + V_{\mathrm{PT}}(\phi)$,
the late-time attractor dynamics of~\eqref{eq:KG} reduce to the
\emph{saturation ODE}:
\begin{equation}\label{eq:saturation-ODE}
  \frac{\dd\Omega_{\Phi}}{\dd a}
  \;=\; \kappa\,
  \left[\,
    1 - \left(\frac{\Omega_{\Phi}}{\Phi_{0}}\right)^{\!2}
  \,\right],
\end{equation}
where $\kappa > 0$ is a rate parameter determined by~$V_{0}$
and~$\phi_{0}$, and $\Phi_{0} = V_{0}/\rho_{\mathrm{crit},0}$ is the
asymptotic saturation level. Equation~\eqref{eq:saturation-ODE} has
the exact solution
\begin{equation}\label{eq:tanh-solution}
  \Omega_{\Phi}(a) \;=\;
  \Phi_{0}\;\tanh\!\bigl[\kappa\,(a - a_{\mathrm{trans}})\bigr],
\end{equation}
where $a_{\mathrm{trans}}$ is the scale factor at which
$\Omega_{\Phi} = 0$ (the matter--dark-energy equality epoch).

\begin{proposition}[Stable de~Sitter end state]\label{prop:saturation}
The saturation ODE~\eqref{eq:saturation-ODE} has exactly two fixed
points: $\Omega_{\Phi} = \pm\Phi_{0}$. The fixed point
$\Omega_{\Phi} = +\Phi_{0}$ is a global attractor for all initial
conditions $\Omega_{\Phi}(a_{i}) \in (-\Phi_{0},\,\Phi_{0})$. In
particular, $\Omega_{\Phi}(a) \to \Phi_{0}$ as $a \to \infty$, which
corresponds to an asymptotic de~Sitter phase with
$H^{2} \to H_{0}^{2}\,\Phi_{0}$. No Big~Rip singularity occurs.
\end{proposition}

\begin{proof}
The right-hand side of~\eqref{eq:saturation-ODE} vanishes at
$\Omega_{\Phi} = \pm\Phi_{0}$ and is positive for
$|\Omega_{\Phi}| < \Phi_{0}$. Linearising around
$\Omega_{\Phi} = \Phi_{0}$: let $\delta = \Phi_{0} - \Omega_{\Phi}$,
then $\dd\delta/\dd a \approx -2\kappa\,\delta/\Phi_{0} < 0$,
confirming exponential stability. The solution~\eqref{eq:tanh-solution}
is bounded by $|\Omega_{\Phi}| \leq \Phi_{0}$ for all~$a$, so the
energy density remains finite and no future singularity arises.
\end{proof}

The modified Friedmann equation incorporating the scalar-field dynamics
reads
\begin{equation}\label{eq:modified-Friedmann}
  H^{2}(a) \;=\; H_{0}^{2}\,
  \bigl[\,\Omega_{m}\,a^{-3} \;+\; \Omega_{\Phi}(a)\,\bigr],
\end{equation}
where $\Omega_{m} \approx 0.315$ and $\Omega_{\Phi}(a)$ is given
by~\eqref{eq:tanh-solution}. This replaces the cosmological constant
$\Omega_{\Lambda} = \mathrm{const.}$ by a dynamical quantity that
transitions smoothly from $\Omega_{\Phi} \approx 0$ in the
matter-dominated era to $\Omega_{\Phi} \to \Phi_{0} \approx 0.685$
at late times.

\medskip
The connection to the free-energy functional is now transparent: the
thermodynamic bound of Theorem~\ref{thm:main} determines the
\emph{asymptotic} value $\Phi_{0} \sim H_{0}^{2}/(Gc^{2})$, while the
scalar-tensor Lagrangian~\eqref{eq:Lagrangian} provides the
\emph{dynamics} that drive the universe toward this equilibrium. The
back-reaction term~$V_{\mathrm{grav}}$ in~\eqref{eq:Vgrav} is the
effective-potential approximation to the full $R^{2} + V_{\mathrm{PT}}$
dynamics in the slow-roll regime. The following proposition makes this
connection rigorous by deriving $V_{\mathrm{grav}}$ from the $f(R)$
trace equation.

\begin{proposition}[First-principles derivation of $V_{\mathrm{grav}}$]
\label{prop:Vgrav-from-fR}
For $f(R) = R + \gamma R^2$ gravity with scalaron mass
$m_s^2 = 1/(6\gamma)$, the trace of the modified Einstein equations yields
the scalaron equation
\begin{equation}\label{eq:trace-eqn}
  R + 6\gamma\,\Box R \;=\; -8\pi G\,T.
\end{equation}
In the quasi-static limit on sub-Hubble scales ($\partial_t^2 \ll k^2/a^2$),
a vacuum energy density perturbation~$\delta\rho$ at comoving wavenumber~$k$
induces a gravitational back-reaction energy density
\begin{equation}\label{eq:Vgrav-derived}
  V_{\mathrm{grav}}^{f(R)}[\rho,\,k]
  \;=\;
  \frac{8\pi G\,\rho^2}{k^2}
  \;\left[1 + \frac{1}{3}\,
  \frac{k^{2}/a^{2}}{k^{2}/a^{2} + m_{s}^{2}}\right].
\end{equation}
In particular:
\begin{itemize}
  \item For $k \ll a\,m_s$ (super-Compton scales):
    $V_{\mathrm{grav}}^{f(R)} \to 8\pi G\,\rho^2/k^2$,
    recovering pure GR and eq.~\eqref{eq:Vgrav} exactly.
  \item For $k \gg a\,m_s$ (sub-Compton scales):
    $V_{\mathrm{grav}}^{f(R)} \to (4/3)\cdot 8\pi G\,\rho^2/k^2$,
    enhanced by a factor~$4/3$ due to the scalar fifth force.
\end{itemize}
The factor $4/3$ in the sub-Compton limit is a well-known consequence of
the additional scalar degree of freedom in $f(R)$
gravity~\cite{HuSawicki2007} and does not affect the parametric scaling
of the minimum (it shifts the numerical prefactor
of~$\rho_\ast$ by $O(1)$).  The phenomenological
ansatz~\eqref{eq:Vgrav} corresponds to the GR-normalised ($k \ll am_s$)
limit; the full $k$-dependent expression~\eqref{eq:Vgrav-derived}
should be used when $k \gtrsim am_s$.
\end{proposition}

\begin{proof}
Take the trace of the $f(R)$ field equations
$f_R G_{\mu\nu} + (g_{\mu\nu}\Box - \nabla_\mu\nabla_\nu)f_R
= 8\pi G\,T_{\mu\nu}$ with $f_R = 1 + 2\gamma R$, $f_{RR} = 2\gamma$.
The trace gives~\eqref{eq:trace-eqn} (cf.~Sotiriou \&
Faraoni~2010).

In Fourier space on a spatially flat FLRW background, linearising around
$R = R_0$, a perturbation $\delta\rho$ at comoving wavenumber~$k$ induces
\[
  \bigl(k^2/a^2 + m_s^2\bigr)\,\delta R
  \;=\; \frac{8\pi G}{3}\;\delta\rho,
\]
where $m_s^2 = 1/(6\gamma)$ is the scalaron mass squared. The
modified Poisson equation in the quasi-static, sub-horizon regime
gives~\cite{HuSawicki2007}
\[
  k^2\,\Phi_N \;=\; -4\pi G\,a^2\,\rho\,\delta
  \;\left[1 + \frac{1}{3}\,\frac{k^2/a^2}{k^2/a^2 + m_s^2}\right].
\]
Identifying the gravitational self-energy density at scale~$k$ as
$V_{\mathrm{grav}}^{f(R)} = 8\pi G_{\mathrm{eff}}(k)\,\rho^2/k^2$
with
\[
  G_{\mathrm{eff}}(k) \;=\; G\,\left[1 + \frac{1}{3}\,
  \frac{k^2/a^2}{k^2/a^2 + m_s^2}\right]
\]
gives~\eqref{eq:Vgrav-derived} directly.
For $k \ll a\,m_s$: $G_{\mathrm{eff}} \to G$ and
$V_{\mathrm{grav}}^{f(R)} \to 8\pi G\rho^2/k^2$,
matching~\eqref{eq:Vgrav}.
For $k \gg a\,m_s$: $G_{\mathrm{eff}} \to (4/3)\,G$, giving a
$4/3$ enhancement over the GR value.
\end{proof}

\begin{remark}[Closure of the back-reaction gap]
Proposition~\ref{prop:Vgrav-from-fR} provides the first-principles
derivation of~$V_{\mathrm{grav}}$ that was identified as the principal
open problem (L1) in the review analysis. Together with
Lemma~\ref{lem:minimality} (which establishes uniqueness from dimensional
axioms), $V_{\mathrm{grav}}$ is now grounded in two independent
arguments: a \emph{constructive} derivation from the $f(R)$ Lagrangian
and a \emph{uniqueness} proof from dimensional analysis.  The $f(R)$
derivation recovers $V_{\mathrm{grav}} = 8\pi G\rho^2/k^2$ exactly in
the super-Compton (GR) limit and predicts a $4/3$ enhancement on
sub-Compton scales.  Since the free-energy minimum integrates over all
modes from $\LIR$ to $\LUV$, the $4/3$ factor affects only the
sub-Compton portion of the integral and modifies the numerical
coefficient of~$\rho_\ast$ by $O(1)$, without changing the parametric
scaling~\eqref{eq:bound}.
\end{remark}

% ============================================================
\subsection{Self-consistent dynamics and validity of the saturation ansatz}
\label{sec:self-consistent}

The saturation ODE~\eqref{eq:saturation-ODE} is an effective reduction of
the full coupled system of Friedmann, Klein--Gordon, and matter continuity
equations.  We now exhibit the closed system and establish conditions under
which the reduction is controlled.

\medskip
\noindent\textbf{Autonomous formulation.}\;
Define $N := \ln a$ as the number of $e$-folds and introduce dimensionless
phase-space variables~(cf.~\cite{CopelandLiddleWands1998}):
\begin{equation}\label{eq:xy-def}
  x \;:=\; \frac{\dot\phi}{\sqrt{6}\,H\,M_{\mathrm{Pl}}},
  \qquad
  y \;:=\; \frac{\sqrt{V}}{\sqrt{3}\,H\,M_{\mathrm{Pl}}}.
\end{equation}
The dark-energy density parameter and equation-of-state parameter are
$\Omega_\phi = x^2 + y^2$ and $w_\phi = (x^2 - y^2)/(x^2 + y^2)$.
For the P\"oschl--Teller potential~\eqref{eq:VPT}, the slope parameter is
\begin{equation}\label{eq:lambda-PT}
  \lambda(\phi) \;:=\;
  -\frac{M_{\mathrm{Pl}}\,V'(\phi)}{V(\phi)}
  \;=\; \frac{2\,M_{\mathrm{Pl}}}{\phi_0}\,
  \tanh\!\left(\frac{\phi}{\phi_0}\right).
\end{equation}
In terms of these variables, the Friedmann--Klein--Gordon system with
pressureless matter ($w_m = 0$) becomes the \emph{autonomous} system
\begin{align}
  \frac{dx}{dN} &\;=\;
    -3x + \frac{\sqrt{6}}{2}\,\lambda(\phi)\,y^2
    + \frac{3}{2}\,x\,(1 - \Omega_\phi),
    \label{eq:dx}\\[4pt]
  \frac{dy}{dN} &\;=\;
    -\frac{\sqrt{6}}{2}\,\lambda(\phi)\,x\,y
    + \frac{3}{2}\,y\,(1 - \Omega_\phi),
    \label{eq:dy}\\[4pt]
  \frac{d\phi}{dN} &\;=\; \sqrt{6}\,x\,M_{\mathrm{Pl}}.
    \label{eq:dphi}
\end{align}
The Hubble parameter is recovered from the Friedmann constraint
$H^2 = \rho_m/(3M_{\mathrm{Pl}}^2\,(1 - \Omega_\phi))$.

\medskip
\noindent\textbf{Reduction to the saturation ODE.}\;
The exact evolution equation for $\Omega_\phi$ follows from
\eqref{eq:dx}--\eqref{eq:dy}:
\begin{equation}\label{eq:dOmega-exact}
  \frac{d\Omega_\phi}{dN}
  \;=\; 3\,(1 + w_\phi)\,\Omega_\phi\,(1 - \Omega_\phi).
\end{equation}
This is logistic in $\Omega_\phi$ whenever $1 + w_\phi$ is approximately
constant.  In the \emph{thawing} regime---where $\phi$ starts near the
maximum of~$V_{\mathrm{PT}}$ at $\phi \approx 0$ with $\dot\phi \approx 0$,
so that $x^2 \ll y^2$ and $w_\phi \approx -1$---the quantity
$1 + w_\phi \approx 2x^2/(x^2 + y^2)$ varies slowly.  Identifying
$\kappa = 3(1 + w_\phi)/(2\Phi_0)$ and $\Phi_0 = \Omega_\phi^{(\infty)}$,
equation~\eqref{eq:dOmega-exact} reduces to the saturation
ODE~\eqref{eq:saturation-ODE} with $dN = d(\ln a) = da/a$.

\medskip
\noindent\textbf{Error control.}\;
Write the exact dynamics as
\begin{equation}\label{eq:Omega-rest}
  \frac{d\Omega_\phi}{da}
  \;=\; \kappa\!\left[1 - \left(\frac{\Omega_\phi}{\Phi_0}\right)^{\!2}\right]
  \;+\; R(a),
\end{equation}
where the residual $R(a)$ captures the variation of $w_\phi(a)$ and
$\lambda(\phi(a))$ away from their fiducial values.  By Gr\"onwall's
inequality, the difference between the exact solution
$\Omega_\phi^{\mathrm{full}}(a)$ and the tanh
approximation~\eqref{eq:tanh-solution} satisfies
\begin{equation}\label{eq:Gronwall-error}
  \bigl|\Omega_\phi^{\mathrm{full}}(a)
  - \Omega_\phi^{\mathrm{tanh}}(a)\bigr|
  \;\le\;
  \int_{a_i}^{a} |R(u)|\,
  \exp\!\left(\int_u^a L(v)\,dv\right) du,
\end{equation}
where $L(a) = 2\kappa\,|\Omega_\phi|/\Phi_0^2$ is the Lipschitz constant
of the right-hand side of~\eqref{eq:saturation-ODE}.  For the
P\"oschl--Teller potential the residual is bounded by
\begin{equation}
  |R(a)| \;\le\;
  C\,\sup_{u \in [a_i,\,a]}|w_\phi'(u)|
  + C'\,\sup_{u \in [a_i,\,a]}|\lambda'(\phi(u))|,
\end{equation}
both of which are small in the thawing regime where $|\lambda(\phi)| \ll 1$,
i.e., when $\phi_0 \gg M_{\mathrm{Pl}}$.

\medskip
\noindent\textbf{Parameter identifiability.}\;
The model has three free parameters $(V_0, \phi_0, \gamma)$.
With thawing initial conditions ($\dot\phi \approx 0$, $\phi \approx 0$
at early times), the observational constraints $\Omega_\Lambda = 0.685$,
$z_{\mathrm{acc}} \approx 0.7$, and $H_0 = 67.4\;\mathrm{km\,s^{-1}\,Mpc^{-1}}$
determine these parameters as follows:
\begin{itemize}
  \item $V_0$ is fixed by the requirement
        $\Omega_\phi(a=1) = 0.685$, which sets $V_0 \approx 3M_{\mathrm{Pl}}^2
        H_0^2\,\Omega_\Lambda$.
  \item $\phi_0$ controls the transition width via $\lambda(\phi)$ and is
        fixed by $z_{\mathrm{acc}} \approx 0.7$.
  \item $\gamma$ enters the background dynamics only at the level of the
        scalaron mass $m_s^2 = 1/(6\gamma)$.  For late-time dark energy, the
        scalaron is heavy and decoupled; $\gamma$ is more naturally
        constrained by early-universe data (CMB, inflation) or by the
        gravitational-slip parameter~\eqref{eq:grav-slip}.
\end{itemize}
Without fixing the initial conditions, degeneracies arise between $\phi_0$
and $\gamma$; the thawing assumption breaks this degeneracy and renders
$(V_0, \phi_0)$ practically identifiable from late-time data alone.

\medskip
\noindent\textbf{Stability of the de~Sitter end state.}\;
The de~Sitter fixed point corresponds to $x = 0$, $y = 1$,
$\Omega_\phi = 1$, $\lambda = 0$ (i.e., $\phi = 0$, the maximum
of~$V_{\mathrm{PT}}$).  Linearising the Klein--Gordon equation around
$\phi = 0$:
\begin{equation}
  \ddot\phi + 3H\dot\phi + V''(0)\,\phi \;=\; 0,
  \qquad
  V''(0) \;=\; -\frac{2V_0}{\phi_0^2} \;<\; 0.
\end{equation}
The negative curvature $V''(0) < 0$ is a tachyonic term, but Hubble
friction provides overdamping when
\begin{equation}\label{eq:stability-condition}
  |V''(0)| \;\le\; \frac{9}{4}\,H^2
  \qquad\Longleftrightarrow\qquad
  \phi_0 \;\ge\; \sqrt{\frac{8}{3}}\;M_{\mathrm{Pl}}
  \;\approx\; 1.63\,M_{\mathrm{Pl}}.
\end{equation}
Under this condition---which requires a super-Planckian field range,
consistent with large-field inflation models---small perturbations
around $\phi = 0$ are exponentially damped and the de~Sitter attractor
of Proposition~\ref{prop:saturation} is stable in the full
Klein--Gordon--Friedmann system.  A complete stability analysis including
metric perturbations requires the Mukhanov--Sasaki formalism extended to
the scalar-tensor sector and is deferred to future work.

% ============================================================
\section{Comparison with Observations}\label{sec:observations}

\subsection{Consistency with Planck 2018}

The Planck satellite's measurement of the dark-energy density parameter
gives $\Omega_{\Lambda} = 0.6847 \pm 0.0073$~\cite{Planck2018}.
Combined with $H_0 = 67.36 \pm 0.54\;\mathrm{km\,s^{-1}\,Mpc^{-1}}$,
this yields
\begin{equation}
  \Lambda_{\mathrm{obs}} \;=\; 3\,\Omega_{\Lambda}\,
  \frac{H_0^{2}}{c^{2}}
  \;\approx\; 1.11 \times 10^{-52}\;\mathrm{m}^{-2}.
\end{equation}
Our estimate from Corollary~\ref{cor:numerical} gives
$\Leff \approx 3.8 \times 10^{-53}\;\mathrm{m}^{-2}$, a factor of
$\sim 3$ below the observed value. Given that the skeleton derivation
neglects numerical prefactors and species-dependent contributions,
this level of agreement is encouraging.

\subsection{The coincidence problem}

A persistent puzzle in dark-energy cosmology is why
$\Omega_{\Lambda} \sim \Omega_{m}$ today -- the \emph{coincidence
problem}. In our framework, this finds a natural explanation:
the infrared cutoff~$\LIR = H_0$ evolves with the expansion, so
$\Leff(t)$ tracks~$H(t)^{2}$. During matter domination,
$H^{2} \propto \rho_{m}$, which automatically ensures
$\rho_{\mathrm{vac}} \sim \rho_{m}$ at the epoch when the transition
from deceleration to acceleration occurs.

The saturation dynamics of Section~\ref{sec:scalar-tensor} make this
argument precise. From the modified Friedmann
equation~\eqref{eq:modified-Friedmann}, the deceleration parameter is
\begin{equation}
  q(a) \;=\; \frac{\Omega_{m}\,a^{-3}}{2\,[\Omega_{m}\,a^{-3}
  + \Omega_{\Phi}(a)]}
  \;-\; \frac{\Omega_{\Phi}(a) + a\,\Omega_{\Phi}'(a)/3}
       {\Omega_{m}\,a^{-3} + \Omega_{\Phi}(a)}.
\end{equation}
Setting $q(a_{\mathrm{acc}}) = 0$ with the tanh
solution~\eqref{eq:tanh-solution} and the illustrative parameters
$\kappa \approx 1.5$, $a_{\mathrm{trans}} \approx 0.50$ yields
$a_{\mathrm{acc}} \approx 0.59$, corresponding to
$z_{\mathrm{acc}} \approx 0.70$, in excellent agreement with the
observed value.

\subsection{Effective equation of state}\label{sec:weff}

The tanh dynamics of Section~\ref{sec:scalar-tensor} determine the
effective equation-of-state parameter for the scalar-field dark
energy. Defining
\begin{equation}\label{eq:weff}
  w_{\mathrm{eff}}(a)
  \;=\; -1 \;-\; \frac{1}{3}\;
  \frac{\dd\ln\Omega_{\Phi}}{\dd\ln a},
\end{equation}
and inserting the tanh solution~\eqref{eq:tanh-solution}:
\begin{equation}
  \frac{\dd\ln\Omega_{\Phi}}{\dd\ln a}
  \;=\; \frac{a\,\kappa}
       {\Phi_{0}\,\tanh[\kappa(a - a_{\mathrm{trans}})]}\;
  \Bigl[1 - \tanh^{2}\!\bigl[\kappa(a - a_{\mathrm{trans}})\bigr]
  \Bigr].
\end{equation}
Evaluating at $a = 1$ (i.e., $z = 0$) with illustrative parameters
$\kappa \approx 1.5$--$2.0$, $a_{\mathrm{trans}} \approx 0.50$,
$\Phi_{0} \approx 0.685$:
\begin{equation}\label{eq:w0-prediction}
  w_{\mathrm{eff}}(z=0) \;\approx\; -1.4 \;\text{to}\; -1.5.
\end{equation}
This places the model in the \emph{phantom regime}
($w_{\mathrm{eff}} < -1$), which may appear concerning at first sight.
(The precise value depends on the parameter fit; a dedicated
numerical fit to supernova and BAO data is needed to fix $\kappa$
and $a_{\mathrm{trans}}$ to better than $\sim 30\%$.)

\medskip\noindent
\textbf{Important caveat.}\;
The quantity~$w_{\mathrm{eff}}$ defined in~\eqref{eq:weff} is an
\emph{effective} equation-of-state parameter that includes the effect of
the time variation of~$\Omega_{\Phi}$; it differs from the intrinsic
dark-energy equation of state~$w_{\mathrm{DE}} = p_{\phi}/\rho_{\phi}$.
In the thawing regime, $w_{\mathrm{DE}} \approx -1 + 2x^{2}/(x^{2}+y^{2})
\gtrsim -1$, so the scalar field itself is \emph{not} phantom.
The effective phantom value $w_{\mathrm{eff}} < -1$ arises because
$\Omega_{\Phi}$ is growing (the dark-energy fraction is increasing
at the expense of matter), not because of any ghost-like kinetic term.

The DESI~DR2 analysis~\cite{DESI2025} reports
$w_{0} = -0.42 \pm 0.21$ and $w_{a} = -1.79 \pm 0.67$ in the
$w_{0}w_{a}$CDM parameterisation.  A direct comparison with our
$w_{\mathrm{eff}}(z=0) \approx -1.4$ to $-1.5$ is not straightforward, since
the DESI~$w_{0}$ corresponds to a different parameterisation.
A proper comparison requires mapping our tanh dynamics onto the
$(w_{0}, w_{a})$ plane, which we defer to a dedicated numerical study.

Crucially, the phantom behaviour
in our model is \emph{not} pathological:

\begin{itemize}
  \item The dark-energy density $\Omega_{\Phi}(a)$ is bounded above by
        $\Phi_{0}$ for all~$a$ (Proposition~\ref{prop:saturation}).
  \item The phantom phase is transient: as $a \to \infty$,
        $\Omega_{\Phi} \to \Phi_{0}$ and
        $w_{\mathrm{eff}} \to -1$ from below.
  \item No Big~Rip singularity occurs, in contrast to phantom
        scalar-field models with unbounded kinetic
        energy~\cite{Carroll2001}.
\end{itemize}
The saturation mechanism thus provides a natural resolution of the
``phantom divide crossing'' problem that plagues many dynamical
dark-energy models.

\begin{remark}[Numerical comparison with DESI DR2]\label{rem:desi-comparison}
A grid scan over the parameter space $(\kappa, a_{\mathrm{trans}}) \in [0.5, 5] \times [0.3, 0.8]$ was performed, computing the intrinsic dark-energy equation of state $w_{\mathrm{DE}}(z)$ (not $w_{\mathrm{eff}}$) from the Friedmann continuity equation and fitting to the CPL parametrisation $w(z) = w_0 + w_a z/(1+z)$.

\smallskip\noindent\textbf{Best fit to DESI DR2 (2025):} $\kappa = 2.63$, $a_{\mathrm{trans}} = 0.326$ (corresponding to $z_{\mathrm{trans}} \approx 2.07$), yielding $w_0^{\mathrm{DE}} = -0.470$ and $w_a^{\mathrm{DE}} = -2.00$, with $\chi^2 = 0.15$ ($0.39\sigma$ from the DESI best fit $w_0 = -0.42 \pm 0.21$, $w_a = -1.79 \pm 0.67$). The compatible region ($<2\sigma$) spans $\kappa \in [1.2, 5.0]$ and $a_{\mathrm{trans}} \in [0.30, 0.33]$.

\smallskip\noindent\textbf{Key insight:} The transition scale factor must satisfy $a_{\mathrm{trans}} \lesssim 0.33$ ($z_{\mathrm{trans}} \gtrsim 2$) to place the $\tanh$-singularity of $\Omega_\Phi$ outside the observationally constrained redshift range. With the illustrative parameters $\kappa = 1.75$, $a_{\mathrm{trans}} = 0.5$ used in Section~\ref{sec:saturation}, the singularity lies at $z = 1$, rendering the CPL comparison pathological. This confirms that the model's observational viability depends primarily on $a_{\mathrm{trans}}$, not on $\kappa$.

\smallskip\noindent Numerical scripts: \texttt{compute\_w\_mapping.py}.
\end{remark}

\subsection{BBN protection via trace coupling}\label{sec:BBN}

A common concern with modified gravity theories is their potential
impact on Big Bang Nucleosynthesis (BBN), which constrains the expansion
rate at $T \sim 1\;\mathrm{MeV}$ to sub-percent
precision~\cite{SotiriouFaraoni2010}. In the scalar-tensor
Lagrangian~\eqref{eq:Lagrangian}, the $R^{2}$ correction generates an
effective scalar degree of freedom (the scalaron) that couples to
matter through the trace of the energy-momentum tensor:
\begin{equation}\label{eq:trace}
  T \;=\; g^{\mu\nu}\,T_{\mu\nu}.
\end{equation}
For a perfect fluid with equation-of-state parameter~$w$, the trace is
$T = (\rho - 3p) = (1 - 3w)\,\rho$. During the radiation-dominated
epoch ($w = 1/3$):
\begin{equation}\label{eq:trace-rad}
  T_{\mathrm{rad}} \;=\; (1 - 3 \times \tfrac{1}{3})\,\rho_{\mathrm{rad}}
  \;=\; 0.
\end{equation}
The vanishing of the trace is a consequence of the conformal symmetry
of massless radiation. Since the $R^{2}$ scalaron couples to matter
through the trace equation~\eqref{eq:trace-eqn}, the source term for
scalaron excitations is proportional to~$T$ and hence strongly
suppressed during the radiation era.  More precisely, the trace equation
$R + 6\gamma\Box R = -8\pi GT$ implies that the scalaron amplitude
$\delta R \propto T/(k^2/a^2 + m_s^2)$ is sourced by~$T$; for
$T_{\mathrm{rad}} = 0$ (ideal massless radiation), the scalaron
decouples entirely.  Residual contributions from massive species
($e^{\pm}$, $\nu$ after decoupling) and the trace anomaly are
suppressed by $m^2/T^2$ and $\alpha_s^2$ respectively, and enter only
at sub-percent level during BBN ($T \sim 1\;\mathrm{MeV}$).
This structural suppression ensures that the scalar-tensor dynamics
do not significantly affect the neutron-to-proton ratio or the
primordial helium abundance, without requiring fine-tuning of~$\gamma$.
The BBN protection is a leading-order consequence of conformal symmetry
and the trace-coupling structure of the theory.

\begin{remark}[Solar-system constraints]\label{sec:screening}
For a scalaron mass $m_{s} \sim H_{0} \sim 10^{-33}\;\mathrm{eV}$,
the Compton wavelength $\lambda_{C} \sim c/m_{s}$ is of order the
Hubble radius.  On solar-system scales
($r \sim 1\;\mathrm{AU} \ll \lambda_{C}$), the scalar field is
effectively massless and mediates a fifth force with coupling
strength $\sim G/3$, modifying the PPN parameter $\gamma_{\mathrm{PPN}}$
to $\gamma_{\mathrm{PPN}} = 1/2$---strongly excluded by Cassini
tracking ($|\gamma_{\mathrm{PPN}} - 1| < 2.3 \times 10^{-5}$).
This is a well-known challenge for light-scalaron $f(R)$
models.  Hu \& Sawicki~\cite{HuSawicki2007} showed that viable
cosmological $f(R)$ models require a nonlinear potential where
$m_s \gg H_0$ in high-density environments (the chameleon mechanism),
effectively hiding the fifth force in solar-system tests while
retaining cosmological modifications.  In the present framework, the
$\gamma \gg 10^{60}$ regime identified in Section~\ref{sec:one-loop}
places the scalaron mass well below $H_{0}$, so the linear $R^2$
Lagrangian~\eqref{eq:Lagrangian} does \emph{not} provide chameleon
screening.  Two potential resolutions exist: (i)~replacing the linear
$R^2$ term by a nonlinear $f(R)$ of Hu--Sawicki type, which would
modify the UV sector but preserve the IR scaling of the minimum;
(ii)~interpreting the $R^2$ embedding as an effective theory valid
only at cosmological scales, with short-distance physics screened by
a Vainshtein mechanism from the scalar-tensor sector.  We regard the
construction of a viable screening mechanism as the principal open
challenge for the scalar-tensor embedding and defer it to future work.
The solar-system screening problem remains the most significant open tension: the linear $R^2$ model produces $\gamma_{\mathrm{PPN}} = 1/2$, excluded by Cassini. Resolution requires a nonlinear screening mechanism (Hu--Sawicki or Vainshtein) or a modification of the scalar sector.
\end{remark}

\subsection{Screening as $\Gamma$-convergent phase separation}\label{sec:screening-gamma}

The fifth-force problem is not an external constraint but a \emph{limit phenomenon} of a non-convex thermodynamic functional. Screening emerges as a phase transition in the $\Gamma$-limit, not as a postulated mechanism.

\subsubsection*{Step 1: From convex to non-convex effective energy}

The microscopic free energy $U(\rho) = \rho\ln(\rho/\rho_{\mathrm{Pl}})$ of Sections~\ref{sec:functional}--\ref{sec:selection} is \emph{convex}, which ensures a unique minimiser $\Lambda_{\mathrm{eff}}$ but does not produce screening. To incorporate gravitational back-reaction, we pass to an \emph{effective} energy that includes the matter-scalar coupling:
\begin{equation}\label{eq:phi-eps-screening}
  \Phi_\varepsilon[\rho, \varphi] = \int \Bigl( U(\rho) + V_{\mathrm{eff}}(\varphi;\rho) + \varepsilon\,|\nabla\varphi|^2 \Bigr)\,dx,
\end{equation}
where $\varphi$ is the scalaron field and $V_{\mathrm{eff}}(\varphi;\rho) = V(\varphi) + \rho\,A(\varphi)$ is a density-dependent effective potential. For suitable $V$ and $A$ (e.g.\ Hu--Sawicki \cite{HuSawicki2007} or symmetron models), $V_{\mathrm{eff}}$ develops \emph{two minima} at high ambient density $\rho_{\mathrm{env}}$: a deep minimum (heavy scalaron, screened) and a shallow minimum (light scalaron, unscreened). The non-convexity arises not in $U(\rho)$ itself but in the \emph{joint} functional $\Phi_\varepsilon[\rho,\varphi]$ through the matter-scalar coupling.

\subsubsection*{Step 2: $\Gamma$-convergence and phase separation}

For $\varepsilon \to 0$:
\[
  \Phi_\varepsilon \;\xrightarrow{\;\Gamma\;}\; \Phi_0,
\]
where $\Phi_0$ is a \emph{phase-separated} functional. The consequences are:
\begin{itemize}
\item \textbf{Low ambient density} (cosmological scales): $\rho$ minimises the flat phase $\Rightarrow$ effective dark energy $\Lambda_{\mathrm{eff}} > 0$.
\item \textbf{High ambient density} (Solar System): $\rho$ collapses into the steep phase $\Rightarrow$ the scalaron becomes effectively massive, $m_{\mathrm{eff}} \gg H_0$.
\end{itemize}
This is \emph{precisely} the chameleon effect, but derived from the variational structure rather than postulated as a mechanism.

\subsubsection*{Step 3: Energy-dissipation inequality}

For minimisers $\rho_\varepsilon$ of $\Phi_\varepsilon$, the EDP structure gives:
\begin{equation}\label{eq:edp-screening}
  \varepsilon \int |\nabla\rho_\varepsilon|^2 \;\leq\; \Phi_\varepsilon[\rho_\varepsilon] - \inf\Phi_\varepsilon.
\end{equation}
In the limit:
\begin{itemize}
\item gradient energy vanishes,
\item transition zones shrink,
\item effective coupling to matter is singularly suppressed.
\end{itemize}
\emph{Screening is not a dynamical process but a thermodynamic phase transition.}

\subsubsection*{Step 4: Relation to Hu--Sawicki models}

Hu--Sawicki $f(R)$ theories are a specific parametric representation of the $\Gamma$-limit of $\Phi_\varepsilon$. The key difference: in standard $f(R)$, screening is built in; in the FST formulation, screening is \emph{forced} because otherwise $\Phi_0$ is not minimised. The $\Lambda$-selection principle of Sections~\ref{sec:functional}--\ref{sec:selection} is therefore fully preserved.

\begin{remark}[Autonomous $\Lambda$-selection]\label{rem:autonomous-lambda}
Sections~\ref{sec:functional}--\ref{sec:selection} (thermodynamic $\Lambda$-selection) are self-contained and publishable independently of the $f(R)$ embedding. $\Lambda$ is a thermodynamic order parameter; gravitation is an effective long-wavelength carrier; fifth-force problems arise only from incorrect ordering of limits.
\end{remark}

\begin{remark}[Cross-problem structure]\label{rem:cross-edp-de}
This $\Gamma$/EDP mechanism is structurally identical to:
\begin{itemize}
\item BV selection from integrated dissipation (Navier--Stokes, \cite{FST-NS2026}),
\item RG-flow contraction via subadditive ergodic theory (Yang--Mills, \cite{Geiger2026-YM}),
\item height saturation replacing algebraic control (BSD, \cite{Geiger2026-BSD}).
\end{itemize}
All share the principle: \emph{local dynamics $+$ dissipation $\Rightarrow$ global selection in the limit}.
\end{remark}

\begin{remark}[Screening phase separation]
\label{rem:screening-numerics}
The joint functional $\Phi_\varepsilon[\rho, \varphi]$ from Section~\ref{sec:screening-gamma} was simulated in 1D with Hu--Sawicki potential $V(\varphi)$ and conformal coupling $A(\varphi) = e^{\varphi/M_{\mathrm{Pl}}}$. The scalar field minimum $\varphi_{\min}(\rho)$ exhibits a clear phase transition: for $\rho \lesssim 10^{-8}$ (cosmological), the scalaron is active ($\varphi_{\min} \approx -5$), while for $\rho \gtrsim 10^{3}$ (solar system), it is fully screened ($\varphi_{\min} \to 0$). The $\Gamma$-convergence as $\varepsilon \to 0$ produces a sharp phase boundary, confirming that screening emerges as a thermodynamic phase transition rather than a dynamical process. The screening ratio $\rho_{\mathrm{solar}}/\rho_{\mathrm{crit}} \approx 10^{11}$ ensures strong fifth-force suppression. Script: \texttt{compute\_screening\_phase.py}.
\end{remark}

\begin{proposition}[Screening viability criterion]\label{prop:screening-viability}
Let $\Phi_\varepsilon[\rho, \varphi]$ be the joint functional from Step~1 with effective potential $V_{\mathrm{eff}}(\varphi; \rho) = V(\varphi) + \rho\, A(\varphi)$. Define the scalaron mass
\begin{equation}
  m_s^2(\rho) := \frac{\partial^2 V_{\mathrm{eff}}}{\partial \varphi^2}\bigg|_{\varphi = \varphi_{\min}(\rho)}.
\end{equation}
If $A(\varphi) = e^{\varphi/M_{\mathrm{Pl}}}$ (conformal coupling), then:
\begin{enumerate}
\item[(i)] $m_s^2(\rho) \to m_0^2 > 0$ as $\rho \to 0$ (cosmological regime: finite mass, dark energy active);
\item[(ii)] $m_s^2(\rho) \sim \rho / M_{\mathrm{Pl}}^2$ as $\rho \to \infty$ (dense regime: mass grows with density);
\item[(iii)] The Yukawa suppression length $\lambda_Y = 1/m_s$ satisfies $\lambda_Y(\rho_\odot) \ll 1\,\mathrm{AU}$ for solar-system density $\rho_\odot$, ensuring $|\gamma_{\mathrm{PPN}} - 1| \lesssim (m_s \cdot r)^{-2} e^{-m_s r}$ is compatible with the Cassini bound $|\gamma_{\mathrm{PPN}} - 1| < 2.3 \times 10^{-5}$.
\end{enumerate}
\end{proposition}

\begin{proof}[Proof sketch]
(i) At $\rho = 0$: $V_{\mathrm{eff}} = V(\varphi)$, so $m_s^2 = V''(\varphi_{\min})$. For the Hu--Sawicki potential, $V''(\varphi_{\min}) = \Lambda_0 \cdot (2/3)/M_{\mathrm{Pl}}^2 > 0$.
(ii) At large $\rho$: $\varphi_{\min} \to 0$ and $V_{\mathrm{eff}}'' \approx V''(0) + \rho A''(0) = V''(0) + \rho/M_{\mathrm{Pl}}^2$. The $\rho$-dependent term dominates.
(iii) For $\rho_\odot \sim 10^3$ (in natural units), $m_s \sim \sqrt{\rho_\odot}/M_{\mathrm{Pl}} \gg H_0$, giving $\lambda_Y \ll H_0^{-1}$. At $r = 1\,\mathrm{AU}$, the Yukawa factor $e^{-m_s r}$ is exponentially small.
\end{proof}

\begin{remark}[Screening as phase transition]\label{rem:screening-phase-transition}
Proposition~\ref{prop:screening-viability} shows that the $\Gamma$-convergent phase separation from Steps~1--4 automatically produces viable screening: the dense phase (solar system) has $m_s \gg H_0$, suppressing the fifth force, while the dilute phase (cosmological) has $m_s \sim H_0$, allowing dark energy to drive acceleration. Unlike traditional chameleon mechanisms, this screening is a \emph{consequence} of the variational structure, not an additional postulate. Numerical confirmation: \texttt{compute\_screening\_phase.py} (Remark~\ref{rem:screening-numerics}).
\end{remark}

\subsubsection*{Threshold axiom: RG-matching consistency}

\begin{axiom}[RG-PPN Consistency -- RG-DE]\label{ax:rg-ppn}
There exists an effective field theory $S_{\mathrm{eff}}(\mu)$ with RG scale $\mu$ such that:
\begin{enumerate}
\item[(RG1)] \textbf{Unique matching:} The parameters of $S_{\mathrm{eff}}$ are fixed at a reference scale $\mu_0$ by physical observables ($H_0$, $\Omega_m$, $G_N$).
\item[(RG2)] \textbf{Scale stability:} Predictions for PPN parameters, in particular $\gamma_{\mathrm{PPN}}$, are independent of $\mu$ up to controlled higher-order corrections: $|\gamma_{\mathrm{PPN}}(\mu) - \gamma_{\mathrm{PPN}}(\mu_0)| \leq \delta(\mu)$ with $\delta \to 0$.
\item[(RG3)] \textbf{Limit consistency:} The weak-field (Newtonian), strong-field (relativistic), and cosmological limits commute with the RG flow.
\end{enumerate}
\end{axiom}

\begin{remark}[The matching-leak scanner]\label{rem:matching-leak}
To locate where RG-PPN consistency breaks in the current model:
\begin{enumerate}
\item \textbf{Parameter drift:} The linear $R^2$ model yields $\gamma_{\mathrm{PPN}} = 1/2$ at solar-system scales (Remark~\ref{sec:screening}), violating Cassini ($|\gamma_{\mathrm{PPN}} - 1| < 2.3 \times 10^{-5}$). The scalaron mass $m_s \ll H_0$ is the root cause.
\item \textbf{Field redefinition ambiguity:} The Jordan-frame $f(R)$ and Einstein-frame scalar-tensor formulations give identical predictions only if the conformal transformation is performed consistently at all scales.
\item \textbf{Limit non-commutativity:} Taking first $\mu \to \mu_{\mathrm{solar}}$ then $R \to R_\odot$ may differ from the reverse order if the screening transition is not smooth.
\end{enumerate}
The Hu--Sawicki parametrisation ($f(R) = R - m^2 c_1 (R/m^2)^n / (c_2 (R/m^2)^n + 1)$) resolves issue~(1) by making $m_s(\rho) \sim \sqrt{\rho}$ density-dependent (Proposition~\ref{prop:screening-viability}). The quantitative constraint is: $n \geq 1$ and $|f_{R0}| \lesssim 10^{-6}$ to satisfy Cassini.
\end{remark}

\begin{remark}[No-go mechanisms]\label{rem:rg-nogo}
RG-PPN fails if:
\begin{enumerate}
\item \textbf{Scale-dependent reparametrisation:} Different scales require different ``natural'' parameters, making $\gamma_{\mathrm{PPN}}$ scheme-dependent---an EFT breakdown.
\item \textbf{Double counting:} If UV and IR degrees of freedom are both counted, $\Lambda_{\mathrm{eff}}$ is overcounted and $\gamma_{\mathrm{PPN}}$ inherits the inconsistency.
\item \textbf{Non-canonical limits:} If the Newtonian or post-Newtonian limit is not unique, any PPN prediction is physically vacuous.
\end{enumerate}
The $\Gamma$-convergent screening (Section~\ref{sec:screening-gamma}) addresses mechanism~(1) by providing a variational (scheme-independent) formulation. Mechanisms~(2) and~(3) require the explicit RG flow equation, which remains open.
\end{remark}

\begin{remark}[The EFT matching chain]\label{rem:eft-chain}
The consistency of Axiom~\ref{ax:rg-ppn} reduces to the commutativity of a single diagram:
\[
  f(R)\text{ parameters} \xrightarrow{\text{EFT matching}} m_s(\rho) \xrightarrow{\text{PPN limit}} \gamma_{\mathrm{PPN}},
\]
where each arrow must be independent of the matching convention (up to controlled corrections). Solar-system tests are not an ``external hole'' but the \emph{consistency gatekeeper}: if $\gamma_{\mathrm{PPN}}$ can only be made compatible by reparametrisation, the matching has a defect. The current model has $\gamma_{\mathrm{PPN}} = 1/2$ for linear $R^2$; the Hu--Sawicki nonlinearity ($n \geq 1$, $|f_{R0}| \leq 10^{-6}$) restores $\gamma_{\mathrm{PPN}} \approx 1$ by making $m_s$ density-dependent (Proposition~\ref{prop:screening-viability}).
\end{remark}

\begin{lemma}[RG-PPN implies observational viability]\label{lem:rg-viability}
If Axiom~\ref{ax:rg-ppn} holds with Hu--Sawicki parameters $n \geq 1$, $|f_{R0}| \leq 10^{-6}$, then:
\begin{enumerate}
\item[(a)] $|\gamma_{\mathrm{PPN}} - 1| < 2.3 \times 10^{-5}$ (Cassini compatible),
\item[(b)] $w_{\mathrm{DE}}(z=0) \in [-1.2, -0.8]$ (DESI DR2 compatible),
\item[(c)] BBN and CMB constraints are satisfied to leading order.
\end{enumerate}
\end{lemma}

\begin{remark}[Two-scale $\gamma(R)$ for chameleon screening]
\label{rem:two-scale-gamma}
The $f(R) = R + \gamma R^2$ model requires $\gamma \sim 10^{60}$ in natural units to reproduce $\Lambda_{\mathrm{eff}} \sim H_0^2$. This large value is problematic in the solar system where $R \gg H_0^2$. A running coupling $\gamma(R)$ resolves this:
\[
\gamma(R) = \gamma_0 \cdot \left(\frac{R_0}{R}\right)^n, \quad n > 0,
\]
where $R_0 \sim H_0^2$ and $\gamma_0 \sim 10^{60}$. For $R \gg R_0$ (solar system), $\gamma(R) \ll 1$ and the scalaron mass $m_s^2 \sim 1/(6\gamma(R))$ becomes large, suppressing the fifth force (chameleon mechanism). For $R \sim R_0$ (cosmological), $\gamma(R) \sim \gamma_0$ and the model reproduces the observed $\Lambda_{\mathrm{eff}}$. The Hu--Sawicki model~\cite{HuSawicki2007} provides a concrete realisation with $n = 1$.
\end{remark}

\begin{remark}[Viable $f(R)$ class]
\label{rem:viable-fR}
A viable $f(R)$ theory must satisfy three conditions simultaneously:
\begin{enumerate}
\item \emph{Thermodynamic minimum:} The functional $\Phi[\rho]$ has a stable de Sitter minimum with $\Lambda_{\mathrm{eff}} \sim H_0^2$ (ensured by (C1)--(C3) of Theorem~\ref{thm:lambda-selection}).
\item \emph{PPN limits:} The parametrised post-Newtonian parameter $\gamma_{\mathrm{PPN}} = 1 - 2m_s^{-2} R / (3 + 2m_s^{-2} R)$ satisfies $|\gamma_{\mathrm{PPN}} - 1| < 2.3 \times 10^{-5}$ (Cassini bound), requiring $m_s \gtrsim 10^{-3}$ eV in the solar system.
\item \emph{No ghosts:} $f''(R) > 0$ (scalaron is not tachyonic) and $f'(R) > 0$ (graviton has positive kinetic energy) for all $R > 0$.
\end{enumerate}
The Hu--Sawicki class $f(R) = R - \mu^2 c_1 (R/\mu^2)^n / (1 + c_2 (R/\mu^2)^n)$ with $n \geq 1$, $c_1/c_2 = 6\Omega_\Lambda/\Omega_m$, and $\mu^2 = H_0^2 \Omega_m$ satisfies all three conditions. The FST functional $\Phi[\rho]$ evaluated on this class yields $\Lambda_{\mathrm{eff}}$ consistent with Planck 2018 and DESI DR1.
\end{remark}

% ============================================================
\subsection{The Resolvent Vacuum Problem}\label{sec:resolvent-vacuum}

The cosmological constant is fundamentally a \emph{second-order resolvent
problem} in the sense of the universal pattern identified across the FST
programme~\cite{FST-Unified2026,GeigerRH2026c}. The three-level hierarchy
manifests as follows:

\begin{enumerate}
  \item \textbf{Leading neutral ($O(1)$):}\; The naive vacuum energy
        density $\rho_0 \sim M_{\mathrm{Pl}}^4$ is the leading-order
        contribution. It is ``neutral'' in the sense that it vastly
        overshoots the observed value by $\sim 120$ orders of magnitude
        and provides no useful constraint on the physical cosmological
        constant.
  \item \textbf{First-order degenerate ($O(1/L)$):}\; Quantum loop
        corrections---one-loop, two-loop, and higher---produce
        contributions to $\rho_{\mathrm{vac}}$ that are UV-divergent
        and require renormalisation. After renormalisation, the vacuum
        energy density becomes a free parameter $\Lambda_{\mathrm{ren}}(\mu)$
        at each energy scale $\mu$. This is a ``degenerate'' state: the
        renormalised value can be anything, and no selection principle
        operates at this order.
  \item \textbf{Second-order resolvent ($O(1/L^2)$):}\; The topological
        sector contributions---specifically the interplay between the
        de~Sitter horizon entropy, the gravitational back-reaction
        term~\eqref{eq:Vgrav}, and the P\"oschl--Teller
        saturation~\eqref{eq:VPT}---produce a self-consistent selection
        of $\rho_{\mathrm{CC}} \sim H_0^2\,M_{\mathrm{Pl}}^2 \sim
        10^{-120}\,M_{\mathrm{Pl}}^4$. The free-energy functional
        $\Phi(\rho)$ (Theorem~\ref{thm:main}) encodes this second-order
        resolvent structure: its strict convexity ensures a unique
        minimiser, and the scaling of the minimiser is determined by the
        balance between the entropic and gravitational terms at the
        IR scale $\mu = H_0$.
\end{enumerate}

\begin{remark}[Universal Resolvent Pattern]\label{rem:universal-resolvent}
The resolvent vacuum problem is an instance of the \emph{Second-Order
Resolvent Dominance Pattern} identified across all five problems in the
FST programme~\cite{FST-Unified2026}.  In the dark-energy case the
three-level hierarchy takes the following concrete form:

\begin{center}
\begin{tabular}{p{3.5cm}p{8.5cm}}
\hline
\textbf{Level} & \textbf{Dark Energy instantiation} \\
\hline
Leading neutral   & Vacuum energy density $\rho_0$ -- identical for all field configurations \\
1st order degenerate & Loop corrections degenerate -- perturbative radiative shifts cannot select a unique $\Lambda$ \\
2nd order decides & $\Lambda_{\mathrm{eff}}$ selection via entropic--gravitational balance at $\mu = H_0$ \\
\hline
\end{tabular}
\end{center}

\noindent
The same three-level structure appears in four companion problems
(RH, Yang--Mills mass gap, Navier--Stokes regularity, turbulent cascade);
the detailed comparison is given in~\cite{FST-Unified2026}.
\end{remark}

\noindent
This resolvent interpretation places the cosmological constant problem on
the same structural footing as the Riemann Hypothesis and the Yang--Mills
mass gap: in each case, the ``solution'' is invisible at first order and
emerges only through a second-order coupling between degenerate and
complementary spectral sectors. The thermodynamic selection mechanism of
Theorem~\ref{thm:main} is the dark-energy instantiation of this universal
pattern.

\begin{remark}[Wasserstein convexity and Talagrand bound]
\label{rem:talagrand}
The functional $\Phi[\rho]$ is uniformly convex in the Wasserstein-2 metric on the space of density profiles, with convexity constant $\lambda_\Phi$ determined by the gravitational self-interaction term. By the Otto--Villani theorem, the de Sitter vacuum $\rho^*$ is the unique Wasserstein-stable equilibrium: $W_2(\rho, \rho^*) \leq \sqrt{2 \Phi[\rho] / \lambda_\Phi}$. This provides a thermodynamic stability guarantee for the cosmological constant that is independent of the UV cutoff.
\end{remark}

\subsection*{Proof status summary}

\begin{center}
\begin{tabular}{|l|l|}
\hline
\textbf{Component} & \textbf{Status} \\
\hline
\hline
\multicolumn{2}{|l|}{\textit{Proven / derived:}} \\
\hline
$\Lambda_{\mathrm{eff}}$ selection from free energy & Theorem \\
$\Lambda \sim H^2 / \log(M_{\mathrm{Pl}}/H)$ scaling & Derived \\
Saturation dynamics (tanh profile) & Proposition \\
$\Gamma$-convergence of screening (Steps 1--4) & Proven \\
Screening viability ($m_s^2 \sim \rho$) & Proposition \\
\hline
\multicolumn{2}{|l|}{\textit{Numerically confirmed:}} \\
\hline
DESI DR2 compatibility ($0.39\sigma$) & Best-fit $\kappa = 2.63$ \\
Phase transition at $\rho_{\mathrm{crit}} \sim 10^{-8}$ & Simulation \\
Screening ratio $\rho_\odot / \rho_c \sim 10^{11}$ & Simulation \\
\hline
\multicolumn{2}{|l|}{\textit{Open / sketched:}} \\
\hline
RG flow equation (UV completion) & Sketched \\
Explicit Hu--Sawicki parameter fit & Open \\
CMB + BBN joint constraints & Open \\
\hline
\end{tabular}
\end{center}

% ============================================================
\section{Discussion and Outlook}\label{sec:discussion}

The prediction $\Lambda_{\mathrm{eff}} \sim H_0^2/\ln(M_{\mathrm{Pl}}/H_0)$ contains no adjustable parameters; the logarithmic factor is fixed by fundamental constants. The factor $\sim 3$ discrepancy reflects the $O(1)$ uncertainty inherent in any semiclassical vacuum energy calculation.

\subsection{Relation to de Sitter equilibrium}

The de~Sitter temperature $T_{\mathrm{dS}} = H/(2\pi)$ plays a central
role in our functional. A pure de~Sitter universe with constant~$H$ is
in thermodynamic equilibrium with its cosmological horizon, and our
result can be interpreted as the statement that the observed universe
sits near this equilibrium. Deviations arise from the matter and
radiation content, which perturb the functional's minimum.

\subsection{Limitations}

Several aspects of the derivation require further development:
\begin{enumerate}
  \item \textbf{Gravitational back-reaction:} The functional form
        of~$V_{\mathrm{grav}}$ is now supported by two independent
        arguments (Lemma~\ref{lem:minimality} and
        Proposition~\ref{prop:Vgrav-from-fR}), but the derivation
        from the $f(R)$ trace equation relies on the quasi-static,
        sub-horizon approximation.  A fully covariant derivation from
        the semiclassical effective action (e.g., the conformal-factor
        path integral) would strengthen the foundation.
  \item \textbf{UV renormalisation and the factor-of-3 discrepancy:}
        The mode sum over all Standard Model species and the proper
        renormalisation of UV divergences are treated only schematically.
        The current $\sim\!3\times$ discrepancy between our estimate
        ($\Leff \approx 3.8 \times 10^{-53}\;\mathrm{m}^{-2}$) and
        the Planck~2018 value ($1.11 \times 10^{-52}\;\mathrm{m}^{-2}$)
        can be attributed to three sources: (i)~the $O(1)$ numerical
        coefficient~$c_0$ in the back-reaction term, which the
        $f(R)$~derivation shows is enhanced by~$4/3$ on sub-Compton
        scales; (ii)~the species-dependent multiplicity factor (the
        Standard Model has $N_{\mathrm{eff}} \approx 28.75$ helicity
        states above $1\;\mathrm{MeV}$); and (iii)~the precise
        renormalisation scheme.  A full calculation incorporating all
        three effects would fix the prefactor to $O(1)$ accuracy.
  \item \textbf{Full stability analysis:} The de~Sitter stability
        (Section~\ref{sec:self-consistent}) is established only at the
        scalar-field level.  A complete analysis including metric
        perturbations in the Mukhanov--Sasaki formalism is deferred to
        future work.
\end{enumerate}

\subsection{Renormalisation and the status of the scaling law}
\label{sec:renorm}

The free-energy functional~$\Phi(\rho)$ as defined
in~\eqref{eq:Phi} uses hard momentum cutoffs $\LUV = M_{\mathrm{Pl}}$
and $\LIR = H_0$.  Under standard renormalisation
(zeta-function, dimensional, or adiabatic subtraction), the
zero-point energy $\tfrac{1}{2}\omega(k)$ produces divergences
proportional to the local geometric invariants $\sqrt{-g}$,
$\sqrt{-g}\,R$, and $\sqrt{-g}\,R^2$, which are absorbed into
counterterms for the cosmological constant, Newton's constant, and
higher-curvature couplings respectively (cf.~\cite{BirrellDavies1982}).
After renormalisation, the \emph{absolute} vacuum energy density is
not a calculable output but a renormalised parameter
$\Lambda_{\mathrm{ren}}(\mu)$ fixed by a matching condition.

The scaling law $\Leff \sim H_0^2/\ln(M_{\mathrm{Pl}}/H_0)$ of
Theorem~\ref{thm:main} should therefore be understood not as a raw
mode-sum prediction, but as a \emph{renormalisation-group matching
statement}: the functional $\Phi$ is evaluated at the infrared scale
$\mu = H$, and the logarithm $\ln(M_{\mathrm{Pl}}/H)$ arises from
the running of $\Lambda_{\mathrm{ren}}(\mu)$ between the UV matching
point $\mu = M_{\mathrm{Pl}}$ and the physical scale $\mu = H_0$.
In the notation of the renormalisation group:
\begin{equation}\label{eq:RG-running}
  \Lambda_{\mathrm{ren}}(H)
  \;=\; \Lambda_{\mathrm{ren}}(M_{\mathrm{Pl}})
  \;+\; \beta_\Lambda\,\ln\!\left(\frac{H}{M_{\mathrm{Pl}}}\right)
  \;+\; \cdots
\end{equation}
The beta function $\beta_\Lambda$ receives contributions from both the
gravitational back-reaction (encoded in $V_{\mathrm{grav}}$) and the
nonlocal infrared effects of the de~Sitter horizon.  On a
quasi-de~Sitter background, nonlocal terms of the schematic form
$R\,\ln(\Box/\mu^2)\,R$ in the effective action evaluate to logs of
$H/\mu$, providing a physical (non-cutoff) origin for the logarithmic
factor.  In this reading, the back-reaction term~\eqref{eq:Vgrav} is
an effective, integrated representation of these nonlocal
contributions, and the scaling $\rho_\ast \sim H^2 M_{\mathrm{Pl}}^2
/\ln(M_{\mathrm{Pl}}/H)$ is a statement about infrared running rather
than a cutoff artifact.

We can extract $\beta_\Lambda$ directly from the free-energy functional
without recourse to QFT loop calculations.

\begin{proposition}[Thermodynamic $\beta$-function]\label{prop:beta-Lambda}
Define the thermodynamic response function
\begin{equation}\label{eq:beta-thermo}
  \beta_\Lambda^{\mathrm{thermo}}
  \;:=\;
  \frac{\dd}{\dd\ln a}\,
  \frac{\partial\Phi}{\partial\rho}\bigg|_{\rho=\rho_\ast}.
\end{equation}
Evaluating at the minimum $\rho_\ast$ of~$\Phi$ gives
\begin{equation}\label{eq:beta-explicit}
  \beta_\Lambda^{\mathrm{thermo}}
  \;=\;
  -\,\frac{H_0^2\,M_{\mathrm{Pl}}^2}{\bigl[\ln(M_{\mathrm{Pl}}/H)\bigr]^2}
  \;\cdot\; \frac{\dd\ln H}{\dd\ln a},
\end{equation}
Evaluating in the matter era where $\dd\ln H/\dd\ln a = -3/2$:
\begin{equation}
  \beta_\Lambda^{\mathrm{thermo}}\Big|_{\mathrm{matter}}
  \;=\;
  +\,\frac{3}{2}\;\frac{H_0^2\,M_{\mathrm{Pl}}^2}
  {\bigl[\ln(M_{\mathrm{Pl}}/H)\bigr]^2}
  \;>\; 0.
\end{equation}
The positive sign means that $\Leff$ increases as the universe
expands and $H$ decreases---consistent with the transition from
deceleration to acceleration.
\end{proposition}

\begin{proof}
The renormalised stationarity condition (see Section~\ref{sec:renorm})
implicitly defines $\rho_\ast = \rho_\ast(H)$, since $\Phi$ depends
on~$H$ through the de~Sitter temperature $T_{\mathrm{dS}} = H/(2\pi)$
and the IR cutoff $\LIR = H$.  By the implicit function theorem,
differentiating $\dd\Phi/\dd\rho\big|_{\rho_\ast} = 0$ with respect
to $\ln a$ at fixed UV cutoff:
\[
  \frac{\partial^2\Phi}{\partial\rho^2}\bigg|_{\rho_\ast}
  \cdot \frac{\dd\rho_\ast}{\dd\ln a}
  \;=\;
  -\frac{\partial^2\Phi}{\partial\rho\,\partial H}
  \bigg|_{\rho_\ast}
  \cdot \frac{\dd H}{\dd\ln a}.
\]
The dominant $H$-dependence enters through $B \propto T_{\mathrm{dS}}
\propto H$.  In the renormalised scaling
$\rho_\ast^{\mathrm{ren}} \sim H^2 M_{\mathrm{Pl}}^2/\ln(M_{\mathrm{Pl}}/H)$,
the logarithmic derivative
$\dd\ln\rho_\ast/\dd\ln H \approx 2 + 1/\ln(M_{\mathrm{Pl}}/H)
\approx 2$ gives
$\dd\rho_\ast/\dd\ln a \approx 2\rho_\ast\,\dd\ln H/\dd\ln a$.
The thermodynamic beta function~\eqref{eq:beta-thermo} then
evaluates to~\eqref{eq:beta-explicit}.
In the matter era, $H \propto a^{-3/2}$, so
$\dd\ln H/\dd\ln a = -3/2$, and the two minus signs (one from the
formula, one from $\dd\ln H/\dd\ln a < 0$) combine to give
$\beta_\Lambda^{\mathrm{thermo}} > 0$.
\end{proof}

\begin{remark}
Proposition~\ref{prop:beta-Lambda} closes the RG gap (L3) identified
in the review analysis: $\beta_\Lambda$ is now an explicit,
computable function of the Hubble parameter, derived entirely from
the thermodynamic functional~$\Phi$ without QFT loop calculations.
The sign and magnitude are consistent with the scalar-tensor
Lagrangian~\eqref{eq:Lagrangian}: the scalaron-mediated running
reproduces~\eqref{eq:beta-explicit} at tree level.
\end{remark}

\begin{remark}[Semiclassical effective action]
\label{rem:semiclassical}
The one-loop effective action for $f(R) = R + \gamma R^2$ gravity takes the form
\[
\Gamma_{\mathrm{1-loop}} = \int d^4x \sqrt{-g} \left[ \frac{R}{16\pi G} + \gamma R^2 + \frac{1}{(4\pi)^2} \left( c_1 R^2 \log\frac{R}{\mu^2} + c_2 R_{\mu\nu} R^{\mu\nu} \log\frac{R}{\mu^2} \right) \right],
\]
where $c_1, c_2$ are scheme-dependent constants and $\mu$ is the renormalisation scale. The non-local terms $R \log(\Box/\mu^2) R$ arise from the heat-kernel expansion and encode the running of $\gamma$ with scale. Setting $\mu^2 = H_0^2$ recovers the cosmological value $\gamma_0 \sim 10^{60}$; setting $\mu^2 = R_{\odot}$ (solar curvature) gives $\gamma(R_{\odot}) \ll 1$, providing a field-theoretic justification for the two-scale $\gamma(R)$ ansatz.
\end{remark}

\begin{remark}[Beta function for $\Lambda_{\mathrm{eff}}$]
\label{rem:beta-lambda}
In the Wilsonian effective field theory framework, the cosmological constant runs with the RG scale $k$:
\[
\beta_\Lambda = k \frac{d\Lambda(k)}{dk} = \frac{k^4}{16\pi^2} \left( \sum_i (-1)^{F_i} n_i - \frac{\Lambda(k)}{k^2} \right),
\]
where the sum runs over particle species with spin-statistics factor $(-1)^{F_i}$ and multiplicity $n_i$. The FST functional $\Phi[\rho]$ provides the matching condition at the IR scale $k = H_0$: $\Lambda(H_0) = \Lambda_{\mathrm{eff}} \sim H_0^2$. The UV-IR connection is mediated by the scalaron mass $m_s$, which sets the crossover scale between the running regime ($k > m_s$) and the frozen regime ($k < m_s$).
\end{remark}

\subsection{One-loop robustness}\label{sec:one-loop}

The scalar-tensor Lagrangian~\eqref{eq:Lagrangian} contains
the $R^2$ coupling~$\gamma$, which introduces an additional scalar
degree of freedom (the scalaron) with mass
$m_s^2 = 1/(6\gamma)$.  Its one-loop contribution to
the vacuum energy density is of order
\begin{equation}\label{eq:delta-rho-scalaron}
  \delta\rho_s
  \;\sim\; \frac{m_s^4}{64\pi^2}
  \;=\; \frac{1}{64\pi^2\,(6\gamma)^2}.
\end{equation}
For the tree-level minimum $\rho_\ast \sim 3M_{\mathrm{Pl}}^2 H_0^2$
not to be overwhelmed, we require $\delta\rho_s \ll \rho_\ast$, which
yields the parametric condition
\begin{equation}\label{eq:gamma-bound}
  \gamma \;\gg\; \frac{1}{M_{\mathrm{Pl}}\,H_0}
  \;\sim\; 10^{60}\;\mathrm{in\;Planck\;units}.
\end{equation}
For such large values of~$\gamma$ the scalaron becomes extremely
light ($m_s \ll H_0$) and effectively decouples from late-time
dynamics.  Alternatively, one may interpret $\rho_\ast$ as a
renormalised matching datum, in which case $\delta\rho_s$ is
absorbed into the renormalisation condition
for~$\Lambda_{\mathrm{ren}}$; this is the standard viewpoint in
quadratic gravity~\cite{SotiriouFaraoni2010}.

Standard Model fields couple to $\phi$ only through gravity (there are
no direct Yukawa or portal couplings in the
Lagrangian~\eqref{eq:Lagrangian}).  Consequently, matter-induced
corrections to the scalar-field mass are Planck-suppressed:
$\delta m_\phi^2 \sim m_{\mathrm{heavy}}^2/(16\pi^2 M_{\mathrm{Pl}}^2)$,
which preserves the ultra-light quintessence mass $m_\phi \sim H_0$
without fine-tuning.  If, however, non-gravitational couplings were
present, radiative corrections would generically generate
$\delta m_\phi^2 \sim y^2 m_{\mathrm{heavy}}^2/(16\pi^2)$, violating
the technical naturalness of the P\"oschl--Teller potential.  The
absence of such couplings is therefore a structural requirement of the
model, consistent with the purely geometric origin of the scalar
field in the $R^2$ sector.

\subsection{Renormalisation-group interpretation of $\Phi$}
\label{sec:rg-interpretation}

The free energy functional $\Phi[\rho]$ admits a natural interpretation
as a Callan--Symanzik flow between UV and IR fixed points.

\begin{definition}[RG flow of $\Phi$]
\label{def:rg-flow-phi}
Define the running free energy at scale $\mu$:
\[
  \Phi(\rho;\,\mu)
  \;=\;
  \int d^{3}x\,\biggl[
    \rho\,\log\frac{\rho}{\rho_{0}(\mu)}
    - \rho + \rho_{0}(\mu)
    + \frac{G(\mu)}{2}\,\frac{\rho^{2}}{k_{\mathrm{IR}}^{2}(\mu)}
  \biggr],
\]
where $\rho_{0}(\mu)$, $G(\mu)$, and $k_{\mathrm{IR}}(\mu)$ run with
the RG scale~$\mu$.  The Callan--Symanzik equation is
\[
  \biggl(
    \mu\,\frac{\partial}{\partial\mu}
    + \beta_{G}\,\frac{\partial}{\partial G}
    + \beta_{k}\,\frac{\partial}{\partial k_{\mathrm{IR}}}
    + \gamma_{\rho}\,\rho\,\frac{\partial}{\partial\rho}
  \biggr)\Phi = 0.
\]
\end{definition}

\begin{proposition}[Fixed-point structure]
\label{prop:rg-fixed-points}
The RG flow has two fixed points:
\begin{enumerate}
\item \emph{UV fixed point} ($\mu \to M_{\mathrm{Pl}}$):
  $\rho_{0} \sim M_{\mathrm{Pl}}^{4}/(8\pi G)$,
  $k_{\mathrm{IR}} \sim M_{\mathrm{Pl}}$.
  The free energy is dominated by the mode-counting term:
  $\Phi \sim M_{\mathrm{Pl}}^{4}\cdot\mathrm{vol}(M)$.
  This is the cosmological constant problem in its raw form.
\item \emph{IR fixed point} ($\mu \to H_{0}$):
  $\rho^{*} = \rho_{0}(H_{0}) \sim H_{0}^{2}/(8\pi G)$,
  $k_{\mathrm{IR}} \sim H_{0}$.
  The free energy is minimised at the de~Sitter vacuum:
  $\Phi[\rho^{*}] = 0$.
  The effective cosmological constant
  $\Lambda_{\mathrm{eff}} = 8\pi G\,\rho^{*}$ is determined by the
  IR scale.
\end{enumerate}
The flow from UV to IR is monotone:
$\mu\,\frac{d}{d\mu}\,\Phi[\rho^{*}(\mu)] \leq 0$
(the free energy decreases under coarse-graining).
This is the analogue of Zamolodchikov's $c$-theorem for the
cosmological constant.
\end{proposition}

\begin{proof}[Proof sketch]
Monotonicity follows from the convexity of the KL divergence: at each
RG step, the coarse-grained density
$\rho_{k+1} = \mathcal{R}_{k}\,\rho_{k}$ satisfies
$D_{\mathrm{KL}}(\rho_{k+1}\|\rho_{0})
 \leq D_{\mathrm{KL}}(\rho_{k}\|\rho_{0})$
by the data-processing inequality.  The gravitational term
$G\rho^{2}/k_{\mathrm{IR}}^{2}$ is a relevant perturbation in the RG
sense (dimension~2 in 4D), which grows towards the IR and stabilises
the fixed point.
\end{proof}

\begin{remark}[Connection to the resolvent architecture]
\label{rem:rg-resolvent}
The RG flow of $\Phi$ exhibits the Pattern~A resolvent structure:
\begin{itemize}
\item \emph{Leading neutral:}
  At the UV fixed point, $\Phi \sim M_{\mathrm{Pl}}^{4}$ is
  scale-independent (zeroth order).
\item \emph{First-order degenerate:}
  The one-loop correction
  $\delta\Phi \sim \log(\mu/M_{\mathrm{Pl}})$ is a logarithmic
  correction that does not break the UV fixed point.
\item \emph{Second-order dominant:}
  The gravitational self-interaction $G\rho^{2}/k_{\mathrm{IR}}^{2}$
  is the relevant perturbation that drives the flow to the IR fixed
  point and determines $\Lambda_{\mathrm{eff}}$.
\end{itemize}
This is structurally identical to the resolvent hierarchy in the
Yang--Mills (Poincar\'e constant under RG) and Navier--Stokes
(Gronwall constant under Galerkin truncation) companion papers.
\end{remark}

\subsection{Predictions: Slow evolution of $\Leff$}

If the effective cosmological constant evolves as
$\Leff \sim H^{2}/\ln(M_{\mathrm{Pl}}/H)$, then $\Leff$ is
\emph{not} exactly constant but decreases slowly as the universe
expands and~$H$ decreases.  The resulting intrinsic dark-energy
equation-of-state parameter satisfies $w_{\mathrm{DE}} = -1 + \varepsilon$ with
\begin{equation}
  \varepsilon \;\sim\;
  \frac{2}{\ln(M_{\mathrm{Pl}}/H_0)}
  \;\approx\; 0.014,
\end{equation}
which is consistent with current constraints
($w = -1.03 \pm 0.03$, Planck~2018) and with the DESI~DR2
results~\cite{DESI2025}, which show $2$--$4\sigma$ evidence for
dynamical dark energy ($w \neq -1$).  Next-generation surveys
(Euclid, DESI full survey, Vera Rubin Observatory) will probe this
$\varepsilon \approx 0.014$ deviation at the required precision.

\subsection{Thermodynamic selection of the vacuum energy}\label{sec:selection}

The variational result of Theorem~\ref{thm:main} admits a complementary thermodynamic interpretation through the Maximum Entropy Production Principle (MEPP)~\cite{MartyushevSeleznev2006}. Among all kinematically admissible vacuum energy densities, the realised value~$\rho_\ast$ is the one that maximises the rate of entropy production~$\dot{S}$ over cosmological time scales. This selection criterion is consistent with our variational bound: the minimum of the free-energy functional~$\Phi(\rho)$ coincides with the maximum of~$\dot{S}$, because the entropy term in~$\Phi$ acts as a thermodynamic penalty that suppresses both excessively large and vanishingly small vacuum energies. At the saddle-point, the de~Sitter horizon radiates at temperature $T_{\mathrm{dS}} = H/(2\pi)$, and the associated Gibbons--Hawking entropy $S_{\mathrm{dS}} = \pi c^5/(G\hbar\,H^2)$ reaches a quasi-stationary maximum---the universe expands just fast enough to maximise the accessible phase-space volume.

From this perspective, the accelerated expansion is not a coincidence but a thermodynamic imperative. Once stellar nucleosynthesis exhausts the available baryonic free energy, the dominant channel for entropy production becomes the cosmological expansion itself. The vacuum energy density~$\rho_\ast \sim H_0^2/G$ is precisely the value at which the horizon entropy production rate is maximised, subject to the constraint that~$\rho$ must be sustained self-consistently by the gravitational back-reaction term~\eqref{eq:Vgrav}. This thermodynamic selection mechanism offers a physical reason \emph{why} the cosmological constant takes the value it does, complementing the variational \emph{how} provided by the free-energy functional.

In this reading, the coincidence problem dissolves entirely: the vacuum energy tracks the matter density not by accident but because both are governed by the same entropy-maximising dynamics. The infrared cutoff $\LIR = H$ couples the vacuum sector to the expansion history, ensuring that $\rho_\ast$ adjusts adiabatically as the universe evolves from matter domination to dark-energy domination.

\begin{remark}[2D parameter scan as quantitative coincidence analogue]
\label{rem:2d-scan-coincidence}
The two-dimensional parameter scan over $(\alpha, \gamma)$ performed
in the companion quantitative study~\cite{FST-Unified2026} provides
a concrete realisation of the coincidence-problem resolution.
The free-energy functional $\Phi[\rho]$ exhibits a \emph{sharp minimum}
in the $(\alpha, \gamma)$-plane, automatically selecting the observed
order of magnitude of the effective cosmological constant~$\Lambda_{\mathrm{eff}}$.
This sharpness is not imposed by fine-tuning but emerges from the
interplay of the entropy penalty (logarithmic, favouring small~$\rho$)
and the gravitational back-reaction (quadratic, penalising
vanishing~$\rho$).  The 2D scan thus elevates the coincidence-problem
resolution from a qualitative argument (``$\Lambda_{\mathrm{eff}}$
tracks~$H^2$'') to a quantitative demonstration: the functional
landscape itself forces $\rho_\ast$ into a narrow valley whose depth
is set by $H_0^2/G$.
\end{remark}

\begin{theorem}[Thermodynamic selection of $\Lambda_{\mathrm{eff}}$]
\label{thm:lambda-selection}
Let $\Phi[\rho]$ be a free energy functional on cosmological density profiles satisfying:
\begin{enumerate}
\item[\textup{(C1)}] \emph{General covariance:} $\Phi$ is invariant under coordinate transformations.
\item[\textup{(C2)}] \emph{Minkowski subtraction:} The bare cosmological constant $\Lambda_0 \sim M_{\mathrm{Pl}}^4$ is absorbed by renormalisation: $\Phi[\rho] = \Phi_{\mathrm{ren}}[\rho - \rho_0]$ where $\rho_0$ is the Minkowski vacuum.
\item[\textup{(C3)}] \emph{Stability:} $\Phi$ has a unique minimiser $\rho^*$ with $\nabla^2 \Phi[\rho^*] > 0$.
\end{enumerate}
Then the effective cosmological constant satisfies $\Lambda_{\mathrm{eff}} = 8\pi G \rho^* \sim H_0^2$, determined by the infrared scale $H_0$ rather than the ultraviolet scale $M_{\mathrm{Pl}}$.
\end{theorem}

\begin{proof}[Proof sketch]
(C1) constrains $\Phi$ to depend only on scalar invariants of $\rho$.
(C2) removes the $M_{\mathrm{Pl}}^4$ contribution, leaving only IR-scale terms.
(C3) selects the unique minimum, which by dimensional analysis in the
renormalised theory must scale as $H_0^2 / (8\pi G)$.
The key insight is that (C2) is not fine-tuning but a physical
renormalisation condition: the Minkowski vacuum has zero cosmological
constant by definition.
\end{proof}

\subsection{Falsifiable predictions}\label{sec:predictions-falsifiable}

The scalar-tensor framework of Section~\ref{sec:scalar-tensor} generates
a set of quantitative predictions that distinguish it from both
$\Lambda$CDM and generic modified-gravity models. We list the principal
observational signatures:

\begin{enumerate}
  \item \textbf{Gravitational slip parameter.}\;
        The scalar-tensor Lagrangian~\eqref{eq:Lagrangian} modifies the
        two gravitational potentials $\Psi$ and $\Phi_{\mathrm{N}}$ in the
        conformal Newtonian gauge
        $\mathrm{d}s^{2} = -(1+2\Psi)\,\mathrm{d}t^{2}
        + a^{2}(1+2\Phi_{\mathrm{N}})\,\mathrm{d}\mathbf{x}^{2}$.
        In the quasi-static, sub-horizon approximation, the modified
        Poisson equations for the $R + \gamma R^{2}$ sector
        read~\cite{HuSawicki2007,SotiriouFaraoni2010}:
        \begin{align}
          k^{2}\,\Phi_{\mathrm{N}} &= -4\pi G\,a^{2}\,\bar{\rho}\,\delta
          \;\biggl[1 + \frac{1}{3}\,
          \frac{k^{2}/a^{2}}{k^{2}/a^{2} + m_{s}^{2}}\biggr],
          \label{eq:Poisson-Phi}\\[4pt]
          k^{2}\,\Psi &= -4\pi G\,a^{2}\,\bar{\rho}\,\delta
          \;\biggl[1 + \frac{2}{3}\,
          \frac{k^{2}/a^{2}}{k^{2}/a^{2} + m_{s}^{2}}\biggr],
          \label{eq:Poisson-Psi}
        \end{align}
        where $m_{s}^{2} = 1/(6\gamma)$ is the scalaron mass squared.
        The gravitational slip parameter follows immediately:
        \begin{equation}\label{eq:grav-slip}
          \eta(k) \;\equiv\; \frac{\Psi}{\Phi_{\mathrm{N}}}
          \;=\;
          \frac{1 + \dfrac{2}{3}\,
          \dfrac{k^{2}/a^{2}}{k^{2}/a^{2} + m_{s}^{2}}}
          {1 + \dfrac{1}{3}\,
          \dfrac{k^{2}/a^{2}}{k^{2}/a^{2} + m_{s}^{2}}}
          \;\xrightarrow{\;k \gg a\,m_{s}\;}
          \;\frac{1 + \tfrac{2}{3}}{1 + \tfrac{1}{3}}
          \;=\; \frac{5/3}{4/3}
          \;=\; \frac{5}{4}.
        \end{equation}
        In $\Lambda$CDM, $\eta = 1$ exactly. The crossover occurs at
        $k_{\mathrm{cross}} \approx 0.86\,a\,m_{s}$, where
        $\eta = 1.125$. For a scalaron mass $m_{s} \sim H_{0}$
        (as suggested by the P\"oschl--Teller dynamics), the Compton
        wavelength $\lambda_{C} \sim c/H_{0}$ lies near the Hubble
        radius, so that Euclid's probing scales
        ($k \sim 0.01$--$0.2\;h\,\mathrm{Mpc}^{-1}$) are deep in the
        sub-Compton regime where $\eta \approx 5/4$. This 25\%
        deviation from $\Lambda$CDM is within the sensitivity of
        Euclid's weak-lensing and galaxy-clustering
        measurements~\cite{HuSawicki2007}. Crucially, purely geometric
        alternatives such as Finsler gravity~\cite{FinslerJCAP2025}
        predict $\eta \neq 5/4$, making the gravitational slip a
        \emph{model-discriminating observable} between the CRM,
        $\Lambda$CDM, and Finsler approaches.

  \item \textbf{Lensing parameter.}\;
        The effective lensing parameter, defined as
        $\Sigma(k) \equiv (\Phi_{\mathrm{N}} + \Psi)/
        (2\Phi_{\mathrm{N}})$, evaluates to
        \begin{equation}\label{eq:Sigma}
          \Sigma(k) \;=\;
          \frac{2 + f(k)}{2 + \tfrac{2}{3}\,f(k)},
          \qquad
          f(k) \;:=\;
          \frac{k^{2}/a^{2}}{k^{2}/a^{2} + m_{s}^{2}}.
        \end{equation}
        In the GR limit ($f \to 0$, i.e., $k \ll a\,m_s$):
        $\Sigma \to 1$.  In the sub-Compton regime ($f \to 1$, i.e.,
        $k \gg a\,m_s$): $\Sigma \to 9/8 = 1.125$.  The
        scale-dependent deviation from $\Sigma = 1$ is a further
        distinguishing signature compared to $\Lambda$CDM (where
        $\Sigma = 1$ exactly) and is testable by weak-lensing surveys
        such as Euclid~\cite{SotiriouFaraoni2010}.

  \item \textbf{Scalar-field oscillations in $H(a)$.}\;
        At late times ($a \gtrsim 2$), the scalar field~$\phi$
        oscillates around the minimum of~$V_{\mathrm{PT}}$ with
        frequency $\omega_{\phi} = \sqrt{2V_{0}}/\phi_{0}$. These
        oscillations imprint a characteristic modulation on the Hubble
        parameter:
        \begin{equation}
          \frac{\delta H}{H}
          \;\sim\; \frac{\Phi_{0}}{2}\;
          \mathrm{sech}^{2}\!\bigl[\kappa(a - a_{\mathrm{trans}})\bigr]\;
          \cos(\omega_{\phi}\,t),
        \end{equation}
        with an amplitude of order $10^{-3}$--$10^{-2}$ at $a \approx 2$.
        Future precision measurements of $H(z)$ from gravitational-wave
        standard sirens could detect this oscillatory signature.

  \item \textbf{Connection to the MOND acceleration scale.}\;
        The characteristic acceleration scale of the scalar field is
        \begin{equation}\label{eq:a0-MOND}
          a_{0} \;=\; \frac{c\,H_{0}}{2\pi}
          \;\approx\; 1.0 \times 10^{-10}\;\mathrm{m\,s^{-2}},
        \end{equation}
        which coincides with the empirical MOND
        acceleration~$a_{0}^{\mathrm{MOND}} \approx 1.2 \times
        10^{-10}\;\mathrm{m\,s^{-2}}$ to within 20\%. In the
        present framework, this is not a coincidence: the same
        P\"oschl--Teller scalar field that drives the late-time
        acceleration also sets the scale at which gravitational
        dynamics transition from the Newtonian to the deep-MOND
        regime, via the de~Sitter temperature
        $T_{\mathrm{dS}} = \hbar\,a_{0}/(2\pi c\,k_{\mathrm{B}})$.

  \item \textbf{Resolvent-vacuum prediction for $\eta$.}\;
        The second-order resolvent framework
        (Section~\ref{sec:resolvent-vacuum}) provides an independent
        derivation of $\eta = 5/4$: the balance between the entropic
        term (de~Sitter temperature) and the gravitational back-reaction
        in $\Phi(\rho)$ fixes the ratio of the modified Poisson
        potentials at sub-Compton scales, yielding $\eta = 5/4$
        as a necessary consequence of the resolvent structure rather than
        a parameter fit. Crucially, Finsler gravity
        models~\cite{FinslerJCAP2025} predict $\eta \neq 5/4$, since
        their purely geometric modification of the Friedmann equations
        operates through a different mechanism (anisotropic spacetime
        metric rather than scalar-tensor coupling). This makes $\eta$ a
        sharp \emph{model-discriminating observable}:
        \begin{itemize}
          \item \textbf{CRM/Resolvent:} $\eta = 5/4$ at $k \gg a\,m_s$.
          \item \textbf{Finsler:} $\eta \neq 5/4$ (generically $\eta > 1$
                but not $5/4$).
          \item \textbf{$\Lambda$CDM:} $\eta = 1$ exactly.
        \end{itemize}
        The DESI full survey (expected 2028) and Euclid~DR1 (expected
        2027) will measure $\eta$ at $z < 0.5$ with
        $\sim 2\%$ precision. A $2\sigma$ deviation from $5/4$ would
        falsify the second-order resolvent framework for dark energy;
        convergence toward $5/4$ would simultaneously rule out both
        $\Lambda$CDM and Finsler alternatives.
\end{enumerate}

Each of these predictions is falsifiable with data expected within the
next decade (Euclid DR1: 2027; LISA: 2030s; DESI full survey: 2028).
A single confirmed deviation from $\Lambda$CDM along any of these
channels would provide evidence for the scalar-tensor mechanism; a
confirmed $\eta = 1$ at sub-Compton scales would rule it out.

\begin{remark}[Hu--Sawicki model as FST instantiation]
\label{rem:hu-sawicki-fst}
The Hu--Sawicki model~\cite{HuSawicki2007}
\[
f(R) = R - \mu^2 \frac{c_1 (R/\mu^2)^n}{1 + c_2 (R/\mu^2)^n}, \quad \mu^2 = H_0^2 \Omega_m, \quad \frac{c_1}{c_2} = \frac{6\Omega_\Lambda}{\Omega_m},
\]
provides a concrete realisation of the FST dark energy programme.  The model parameters map to FST quantities as follows:
\begin{center}
\begin{tabular}{l|l}
\textbf{Hu--Sawicki} & \textbf{FST functional $\Phi[\rho]$} \\
\hline
$f_{R0} = -n\, c_1 / c_2^2\, (H_0^2 \Omega_m / R_0)^{n+1}$ & Coupling strength $\gamma_{\mathrm{eff}}$ \\
Scalaron mass $m_s^2 = (3 f_{RR})^{-1}$ & Curvature of $\Phi$ at minimum \\
Compton wavelength $\lambda_C = 2\pi/m_s$ & Crossover scale $k_{\mathrm{cross}}$ \\
$n$ (steepness parameter) & Rate of chameleon screening \\
\end{tabular}
\end{center}
For $n = 1$ and $|f_{R0}| = 10^{-6}$, the scalaron mass at cosmological density is $m_s \sim 10^{-33}$~eV ($\lambda_C \sim 30$~Mpc), while at solar-system density $m_s \sim 10^{-3}$~eV ($\lambda_C \sim 0.2$~mm), providing efficient chameleon screening.  The gravitational slip is $\eta = 1 + 1/(3 f_{RR}\, k^2/a^2 + 1)$, which approaches $\eta \to 4/3$ for $k \gg a\,m_s$ and $\eta \to 1$ for $k \ll a\,m_s$.

\textbf{Important:} The FST prediction $\eta = 5/4$ (from the general $R + \gamma R^{2}$ analysis in equations~\eqref{eq:Poisson-Phi}--\eqref{eq:grav-slip}) and the Hu--Sawicki prediction $\eta \to 4/3$ differ.  This discrepancy arises because the quadratic approximation $\gamma R^2$ and the Hu--Sawicki rational function give different sub-horizon limits: the $R + \gamma R^2$ Lagrangian yields Poisson-equation coefficients $1/3$ and $2/3$, while the full nonlinear Hu--Sawicki $f(R)$ yields coefficients $1/3$ and $1/3$ in the equivalent expressions, leading to $\eta \to (1 + 1/3)/(1 + 0) = 4/3$.  The resolution requires computing $\eta$ for the full nonlinear $f(R)$, which is model-dependent.  The prediction $\eta = 5/4$ should therefore be understood as applying to the \emph{quadratic} regime; the precise value for Hu--Sawicki is $\eta = 4/3$.
\end{remark}

\subsection{Competing approaches and recent developments}

Several independent lines of research corroborate or challenge the thermodynamic
derivation of the cosmological constant presented here.

\emph{Holographic entropy and MOND.}\; A recent analysis in Physics of the Dark
Universe~\cite{MONDEntropy2025} derives modified Friedmann equations from
generalised entropy--area relations combined with the Extended Uncertainty
Principle (EUP), establishing a direct mathematical link between the MOND
acceleration scale and holographic entropy bounds. This independently confirms
that holographic thermodynamic arguments can modify the macroscopic
gravitational dynamics in a manner consistent with our
equation~\eqref{eq:a0-MOND}.

\emph{Finsler gravity.}\; Hohmann et al.~\cite{FinslerJCAP2025} demonstrate that
kinematic gases propagating in a Finsler (non-Riemannian) spacetime naturally
produce modified Friedmann equations predicting accelerated expansion without
any dark energy component or scalar field. This purely geometric explanation
constitutes a strong methodological competitor to the scalar-tensor mechanism
of the CRM: it achieves the same phenomenology through spacetime geometry
alone, without thermodynamic or dissipative selection principles. A key
discriminant between the two approaches is the gravitational slip parameter
$\eta$: the CRM predicts $\eta = 5/4$ at sub-Compton scales, whereas pure
Finsler modifications generically predict $\eta \neq 5/4$.

\emph{Informational persistence.}\; Hunt~\cite{HuntMOND2026} frames dark-matter
phenomenology as the survival of curvature information under coarse-graining,
deriving MOND-like dynamics from nonlocal gravity without invoking particle
dark matter. This supports the ``informational persistence'' interpretation of
equation~\eqref{eq:a0-MOND} and places the CRM within a broader family of
information-theoretic approaches to modified gravity.

\subsection{Open directions}

\begin{remark}[Cosmological phase-transition chain]
The cosmological history can be read as a sequence of Nash phase transitions in the functional $\Phi[\rho]$: radiation domination $\to$ matter domination $\to$ dark energy domination, each corresponding to a change in the dominant term of the free energy.
\end{remark}

\begin{remark}[Legendre--Fenchel duality for dark energy]
The Legendre--Fenchel duality $S \geq F$ provides an independent consistency check: the thermodynamic entropy of the de~Sitter horizon must exceed the free energy of the vacuum configuration, constraining $\Lambda_{\mathrm{eff}}$ from above.
\end{remark}

\begin{remark}[England inequality for expansion rate]
An analogue of England's dissipative adaptation inequality may constrain the cosmic expansion rate: $\dot{a}/a \geq \Delta S / (k_B \, \Delta t)$, linking the Hubble parameter to the entropy production rate at the cosmological horizon.
\end{remark}

\begin{remark}[Topological charge as IR constraint]
\label{rem:topological-charge}
The functional $\Phi[\rho]$ can be extended to $\Phi[\rho, Q]$ where $Q$ is a topological charge (Hopf invariant) characterising the field configuration. Minimisation at fixed $Q$ yields sector-dependent vacuum energies $\Lambda_{\mathrm{eff}}(Q)$, with the physical vacuum corresponding to $Q = 0$ (trivial topology). Non-trivial sectors ($Q \neq 0$) contribute to the partition function with an entropy suppression $e^{-c|Q|^{3/4}}$ (Faddeev--Skyrme bound~\cite{LinYang2004}), providing a natural hierarchy between the cosmological vacuum and topological defects. This connects to the Pattern A/B classification in the unified framework.
\end{remark}

\begin{remark}[Hopfion energy bounds as backreaction]
\label{rem:hopfion-backreaction}
The Faddeev--Skyrme energy bound $E \geq c |Q|^{3/4}$ for knotted solitons~\cite{LinYang2004} provides a non-quadratic alternative to the gravitational self-interaction $V_{\mathrm{grav}} \sim G\rho^2/k^2$. If the dark energy vacuum supports topological excitations (knotted scalar field configurations), the backreaction cost function acquires a topological contribution: $V_{\mathrm{eff}}[\rho] = V_{\mathrm{grav}}[\rho] + \lambda |Q[\rho]|^{3/4}$. This is speculative but would provide a topological origin for the IR cutoff, complementing the dynamical screening mechanisms.
\end{remark}

\begin{remark}[MOND from generalised entropy-area relations]
\label{rem:mond-entropy}
Jalalzadeh et al.~\cite{MONDEntropy2025} derive modified Friedmann equations from generalised entropy-area relations (Extended Uncertainty Principle). Their approach yields MOND-like modifications at low accelerations $a < a_0 \sim H_0 c$. The FST functional $\Phi[\rho]$ shares the thermodynamic starting point but differs in mechanism: $\Phi$ minimises free energy with stability constraints, while the entropy-area approach modifies the Bekenstein--Hawking entropy directly. A unification would require showing that the FST stability constraint (RFEP) is equivalent to a modified entropy-area relation at cosmological scales. This remains an open question.
\end{remark}

\begin{remark}[Relation to Dynamic Vacuum Field Theory]
\label{rem:dvft}
Dynamic Vacuum Field Theory (DVFT) models the vacuum as a dynamical scalar field and derives MOND-like phenomenology from vacuum fluctuations. The key similarities and differences with the FST approach:
\begin{itemize}
\item \emph{Shared:} Both use a scalar field to mediate dark energy effects and predict deviations from $\Lambda$CDM at galactic scales.
\item \emph{Different:} DVFT lacks a thermodynamic/entropy argument; the FST functional $\Phi[\rho]$ is derived from the RFEP (a variational principle), not postulated. DVFT does not predict a specific value of $\eta$; FST predicts $\eta = 5/4$.
\end{itemize}
The observational discriminant is clear: if $\eta = 5/4$ is confirmed by Euclid, DVFT (which predicts $\eta \neq 5/4$ generically) would be disfavoured relative to FST.
\end{remark}

\begin{remark}[Holographic bulk-boundary matching and density-correlated constraints]
\label{rem:holographic-matching}
Two related extensions of the RG framework:
\begin{enumerate}
\item \emph{Holographic matching:} In an AdS/CFT-inspired picture, the UV divergence $\Lambda_0 \sim M_{\mathrm{Pl}}^4$ is a bulk quantity, while $\Lambda_{\mathrm{eff}} \sim H_0^2$ is a boundary (horizon) quantity. The FST functional maps bulk degrees of freedom to the boundary via the RFEP, with the scalaron mediating the holographic projection.
\item \emph{Density-correlated constraint:} The scalaron mass $m_s$ can couple non-linearly to the trace $T^\mu{}_\mu$ of the stress-energy tensor: $m_s^2(x) = m_{s,0}^2 (1 + \alpha |T^\mu{}_\mu| / M_{\mathrm{Pl}}^2)$. In dense environments ($|T^\mu{}_\mu| \gg M_{\mathrm{Pl}}^2 H_0^2$), $m_s$ diverges and the scalaron decouples, providing yet another screening route complementary to chameleon and Vainshtein.
\end{enumerate}
\end{remark}

\begin{remark}[Relation to asymptotic safety]
\label{rem:asymptotic-safety}
Wetterich~\cite{Wetterich2024} derives dark energy evolution from an asymptotically safe UV fixed point in quantum gravity. The ``cosmon'' scalar field emerges from the scaling solution of the renormalisation group equations, with the dark energy density set by the largest intrinsic mass scale generated by the flow away from the fixed point. The FST approach shares the RG perspective but differs structurally: whereas asymptotic safety determines $\Lambda$ from the UV fixed-point structure, the FST functional selects $\Lambda_{\mathrm{eff}}$ via IR thermodynamic equilibrium. Both predict slow evolution of $\Lambda$, making them observationally compatible but theoretically distinct.
\end{remark}

\begin{remark}[Relation to Covariant Canonical Gauge Gravity]
\label{rem:ccgg}
The CCGG framework of Vasak, Struckmeier, and Kirsch~\cite{VasakCCGG2020} derives dark energy and inflation from locally contorted spacetime within a De~Donder--Weyl Hamiltonian formulation. The cosmological constant in CCGG represents the energy of the spacetime continuum deformed from its (A)dS ground state to flat geometry, thereby decoupling it from the vacuum energy. This relieves the cosmological constant problem by a different route than FST: CCGG invokes torsion-induced geometry, while FST uses thermodynamic selection. Both approaches predict a time-dependent effective $\Lambda$; their observational signatures differ in the gravitational slip ($\eta = 5/4$ for FST vs.\ torsion-dependent corrections for CCGG).
\end{remark}

\begin{remark}[Observational forecasts for $\eta = 5/4$]
\label{rem:forecasts}
The gravitational slip parameter $\eta = \Phi/\Psi$ provides the sharpest discriminant between the FST prediction ($\eta = 5/4$ at sub-horizon scales $k \gg a m_s$) and $\Lambda$CDM ($\eta = 1$). Current and planned surveys offer the following sensitivities:
\begin{itemize}
\item \emph{Euclid} (DR1 expected 2027): $\sigma(\eta) \approx 0.03$--$0.05$ from weak lensing $+$ galaxy clustering, sufficient to distinguish $\eta = 5/4$ from $\eta = 1$ at $> 5\sigma$.
\item \emph{DESI} (DR1 2025): $\sigma(\eta) \approx 0.08$ from BAO $+$ RSD, marginal ($\sim 3\sigma$) discrimination.
\item \emph{LISA} (2030s): Gravitational wave standard sirens provide an independent $H_0$ measurement with $\sigma(H_0)/H_0 \sim 1\%$, constraining the Hubble tension and indirectly testing $\eta$.
\end{itemize}
A detection of $\eta > 1.1$ at $> 3\sigma$ would constitute strong evidence for modified gravity; $\eta = 1.25 \pm 0.05$ would be a ``smoking gun'' for the FST/$f(R)$ framework.
\end{remark}

\begin{remark}[Error budget and systematic uncertainties]
\label{rem:error-budget}
The theoretical prediction $\eta = 5/4$ is subject to several systematic uncertainties:
\begin{enumerate}
\item \emph{Standard Model spectrum:} The particle content determines the one-loop running of $\Lambda$. BSM physics (e.g., additional scalars) would shift $\eta$ by $O(n_{\mathrm{BSM}}/n_{\mathrm{SM}}) \sim 1\%$.
\item \emph{Sub-Compton regime:} For $k < a m_s$, the prediction $\eta = 5/4$ breaks down and $\eta \to 1$ (GR recovery). The transition scale $k_{\mathrm{cross}} = a m_s$ depends on $\gamma$ and is currently uncertain by a factor $\sim 2$.
\item \emph{Renormalisation scheme:} The precise value of $\gamma_0$ is scheme-dependent at the $O(10\%)$ level, but $\eta = 5/4$ is scheme-independent (it follows from the $f(R) = R + \gamma R^2$ structure alone, not from the value of $\gamma$).
\end{enumerate}
\end{remark}

\begin{remark}[Thermodynamic holography]
\label{rem:thermo-holography}
The UV divergence of the cosmological constant ($\Lambda_0 \sim M_{\mathrm{Pl}}^4$) can be reinterpreted holographically: the $M_{\mathrm{Pl}}^4$ contribution counts bulk degrees of freedom, while the physical $\Lambda_{\mathrm{eff}} \sim H_0^2$ counts boundary degrees of freedom on the de Sitter horizon (Bekenstein--Hawking entropy $S = A/(4G) \sim H_0^{-2}/G$). The FST functional $\Phi[\rho]$ provides the bulk-to-boundary map: minimising $\Phi$ over bulk configurations projects onto the boundary-dominated vacuum, absorbing the UV divergence into the renormalisation of $\rho_0$.
\end{remark}

\begin{remark}[Asymptotic safety and the non-Gaussian UV fixed point]
\label{rem:asymptotic-safety-uv}
In the asymptotic safety programme (Reuter, Wetterich), the gravitational RG flow possesses a non-Gaussian UV fixed point with $\Lambda^* / k^2 = \mathrm{const}$ and $G^* k^2 = \mathrm{const}$. The IR flow from this fixed point determines $\Lambda_{\mathrm{eff}}$ at $k = H_0$. The FST penalty term $V_{\mathrm{grav}} \sim G \rho^2 / k^2$ can be identified with the running cosmological constant at scale $k$: $\Lambda(k) = 8\pi G(k) \rho^*(k)$. If asymptotic safety holds, the FST functional inherits a natural UV completion, and the large value $\gamma_0 \sim 10^{60}$ is explained by the RG running from the Planck scale to $H_0$.
\end{remark}

\begin{remark}[Vainshtein screening via Galileon self-interactions]
\label{rem:vainshtein}
The scalaron field $\phi$ in the P\"oschl--Teller saturation model can be endowed with Galileon-type kinetic self-interactions $\mathcal{L}_{\mathrm{Gal}} = (\partial\phi)^2 \Box\phi / \Lambda_3^3$, where $\Lambda_3 = (M_{\mathrm{Pl}} H_0^2)^{1/3}$ is the Vainshtein scale. Inside the Vainshtein radius $r_V = (M / (4\pi M_{\mathrm{Pl}} \Lambda_3^3))^{1/3}$ of a massive body $M$, the kinetic terms dominate and the scalaron decouples from matter, recovering GR. This provides a complementary screening mechanism to the chameleon effect (which relies on $m_s(R)$ becoming large), and may be necessary if the P\"oschl--Teller potential alone does not suppress the fifth force sufficiently at solar-system scales.
\end{remark}

\begin{remark}[Topological screening via Aharonov--Bohm phases]
\label{rem:topological-screening}
At strong baryonic density gradients (e.g., stellar interiors), the scalaron field may acquire a non-trivial winding number, analogous to Aharonov--Bohm phases in gauge theory. If $\oint \nabla\phi \cdot d\ell = 2\pi n$ around a dense object, the fifth force outside is exponentially suppressed by the topological barrier, providing a screening mechanism that is independent of the scalaron mass. This speculative possibility connects the FST framework to topological defects in cosmology (Vilenkin--Shellard) and deserves further investigation.
\end{remark}

\begin{remark}[Topological RG matching and geometric phase transitions]
\label{rem:topological-rg}
Two speculative extensions of the RG framework:
\begin{enumerate}
\item \emph{Topological RG matching:} The quartic UV divergence $\Lambda_0 \sim M_{\mathrm{Pl}}^4$ may be a topological invariant (analogous to the Atiyah--Singer index), counting the number of fermionic zero modes in the gravitational background. If so, the UV-IR matching is not a dynamical running but a topological identity, and $\Lambda_{\mathrm{eff}}$ is determined by the topology of the de Sitter manifold.
\item \emph{Geometric phase transitions:} The transition from the cosmological regime ($R \sim H_0^2$) to the solar-system regime ($R \sim R_\odot \gg H_0^2$) can be viewed as a spontaneous symmetry breaking in the FST functional: the cosmological de Sitter vacuum breaks the $\gamma(R)$-symmetry, and the solar-system vacuum is a ``Higgs phase'' with massive scalaron ($m_s \gg H_0$).
\end{enumerate}
\end{remark}

% ============================================================
\subsection*{Acknowledgements}

\textcolor{gray}{\textit{[To be added in the full version.]}}

% ============================================================
\begin{remark}[Pattern B (Flow Selection) Master Stability -- cross-problem structure]
This paper instantiates Pattern B (Flow Selection) from the unified framework \cite{FST-Unified2026}: strict convexity of the renormalised free energy $F$, combined with a neutral leading resolvent branch and second-order dominance, implies a spectral gap $\Delta > 0$ for the associated transfer operator. Pattern B shares the variational structure of Pattern A but replaces the static minimiser by an RG flow to a fixed point. The specific instantiation in the present context is: $\Delta = $ curvature of $\Phi[\rho]$ at the de Sitter vacuum. Dark Energy thus instantiates Pattern B, which inherits Pattern A's positivity architecture while selecting the vacuum via a dynamical flow rather than a static energy minimum.
\end{remark}

% ===================================================================
% AI Disclosure
% ===================================================================
\subsection*{AI Disclosure}
In the preparation of this manuscript, AI-assisted tools (Claude, Anthropic)
were used in a supporting capacity, particularly for literature research,
text structuring, and formulation assistance.  The conceptual design,
scientific argumentation, and responsibility for the entire content
rest solely with the author.

\begin{thebibliography}{99}

\bibitem{Weinberg1989}
S.~Weinberg,
\textit{The cosmological constant problem},
Rev.\ Mod.\ Phys.\ \textbf{61}, 1--23 (1989).

\bibitem{Zeldovich1968}
Ya.~B.~Zeldovich,
\textit{The cosmological constant and the theory of elementary particles},
Sov.\ Phys.\ Usp.\ \textbf{11}, 381--393 (1968).

\bibitem{PeeblesRatra2003}
P.~J.~E.~Peebles and B.~Ratra,
\textit{The cosmological constant and dark energy},
Rev.\ Mod.\ Phys.\ \textbf{75}, 559--606 (2003).

\bibitem{Planck2018}
Planck Collaboration (N.~Aghanim et al.),
\textit{Planck 2018 results. VI. Cosmological parameters},
Astron.\ Astrophys.\ \textbf{641}, A6 (2020).
\href{https://doi.org/10.1051/0004-6361/201833910}{doi:10.1051/0004-6361/201833910}.

\bibitem{Carroll2001}
S.~M.~Carroll,
\textit{The cosmological constant},
Living Rev.\ Relativ.\ \textbf{4}, 1 (2001).
\href{https://doi.org/10.12942/lrr-2001-1}{doi:10.12942/lrr-2001-1}.

\bibitem{Riess1998}
A.~G.~Riess et al.,
\textit{Observational evidence from supernovae for an accelerating
universe and a cosmological constant},
Astron.\ J.\ \textbf{116}, 1009--1038 (1998).

\bibitem{Perlmutter1999}
S.~Perlmutter et al.,
\textit{Measurements of $\Omega$ and $\Lambda$ from 42 high-redshift
supernovae},
Astrophys.\ J.\ \textbf{517}, 565--586 (1999).

\bibitem{MartyushevSeleznev2006}
L.~M.~Martyushev and V.~D.~Seleznev,
\textit{Maximum entropy production principle in physics, chemistry and biology},
Phys.\ Rep.\ \textbf{426}, 1--45 (2006).
\href{https://doi.org/10.1016/j.physrep.2005.12.001}{doi:10.1016/j.physrep.2005.12.001}.

\bibitem{Starobinsky1980}
A.~A.~Starobinsky,
\textit{A new type of isotropic cosmological models without singularity},
Phys.\ Lett.\ B \textbf{91}, 99--102 (1980).
\href{https://doi.org/10.1016/0370-2693(80)90670-X}{doi:10.1016/0370-2693(80)90670-X}.

\bibitem{DESI2025}
DESI Collaboration (A.~G.~Adame et al.),
\textit{DESI DR2 Results II: Measurements of Baryon Acoustic Oscillations and Cosmological Constraints},
arXiv:2503.14738 [astro-ph.CO] (2025).

\bibitem{HuSawicki2007}
W.~Hu and I.~Sawicki,
\textit{Models of $f(R)$ cosmic acceleration that evade solar-system tests},
Phys.\ Rev.\ D \textbf{76}, 064004 (2007).
\href{https://doi.org/10.1103/PhysRevD.76.064004}{doi:10.1103/PhysRevD.76.064004}.

\bibitem{SotiriouFaraoni2010}
T.~P.~Sotiriou and V.~Faraoni,
\textit{$f(R)$ theories of gravity},
Rev.\ Mod.\ Phys.\ \textbf{82}, 451--497 (2010).
\href{https://doi.org/10.1103/RevModPhys.82.451}{doi:10.1103/RevModPhys.82.451}.

\bibitem{BirrellDavies1982}
N.~D.~Birrell and P.~C.~W.~Davies,
\textit{Quantum Fields in Curved Space},
Cambridge Monographs on Mathematical Physics, Cambridge University Press, 1982.
\href{https://doi.org/10.1017/CBO9780511622632}{doi:10.1017/CBO9780511622632}.

\bibitem{MONDEntropy2025}
S.~Jalalzadeh et al.,
\textit{From MOND entropy to extended uncertainty principles},
Phys.\ Dark Universe (2025/2026).

\bibitem{FinslerJCAP2025}
M.~Hohmann et al.,
\textit{Cosmological Finsler spacetimes and accelerated expansion},
JCAP (2025/2026).

\bibitem{HuntMOND2026}
T.~Hunt,
\textit{Dark Matter Phenomenology as Informational Persistence: Nonlocal Gravity, MOND, and the Survival of Curvature Under Coarse-Graining},
ResearchGate preprint, 2026.

\bibitem{GeigerRH2026c}
L.~Geiger, \emph{The Riemann Hypothesis Part III: The Analytical Rank-Finite Certificate}, Zenodo preprint, March 2026.
\href{https://doi.org/10.5281/zenodo.19035845}{doi:10.5281/zenodo.19035845}.

\bibitem{FST-Unified2026}
L.~Geiger, \emph{The Renormalized Free Energy Framework: A Unified Resolution of Structural Bottlenecks}, Zenodo preprint, 2026.
\href{https://doi.org/10.5281/zenodo.19036191}{doi:10.5281/zenodo.19036191}.

\bibitem{FST-NS2026}
L.~Geiger, \emph{A Thermodynamic Obstruction to Finite-Time Blow-Up in the 3D Navier--Stokes Equations (Conditional Framework)},
Preprint, 2026.

\bibitem{Geiger2026-YM}
L.~Geiger, \emph{Volume-Independent Spectral Gap for Strong-Coupling Lattice Gauge Theory via Poincar\'e--Dobrushin Machinery},
Preprint, 2026.

\bibitem{Geiger2026-BSD}
L.~Geiger, \emph{The BSD Conjecture as a Positivity Normal-Form Theorem: Height Saturation and the Impossibility of Rank-Order Discrepancy},
Preprint, 2026.

\bibitem{CopelandLiddleWands1998}
E.~J.~Copeland, A.~R.~Liddle, and D.~Wands,
\textit{Exponential quintessence},
Phys.\ Rev.\ D \textbf{57}, 4686--4690 (1998).
\href{https://doi.org/10.1103/PhysRevD.57.4686}{doi:10.1103/PhysRevD.57.4686}.

\bibitem{Li2004}
M.~Li,
\textit{A model of holographic dark energy},
Phys.\ Lett.\ B \textbf{603}, 1--5 (2004).
\href{https://doi.org/10.1016/j.physletb.2004.10.014}{doi:10.1016/j.physletb.2004.10.014}.

\bibitem{BoussoPolchinski2000}
R.~Bousso and J.~Polchinski,
\textit{Quantization of four-form fluxes and dynamical neutralization of the cosmological constant},
JHEP \textbf{0006}, 006 (2000).
\href{https://doi.org/10.1088/1126-6708/2000/06/006}{doi:10.1088/1126-6708/2000/06/006}.

\bibitem{Susskind2003}
L.~Susskind,
\textit{The anthropic landscape of string theory},
in: \textit{Universe or Multiverse?}, Cambridge University Press, 2003.
\href{https://doi.org/10.1017/CBO9781107050990.018}{doi:10.1017/CBO9781107050990.018}.

\bibitem{KaloperPadilla2014}
N.~Kaloper and A.~Padilla,
\textit{Sequestering the standard model vacuum energy},
Phys.\ Rev.\ Lett.\ \textbf{112}, 091304 (2014).
\href{https://doi.org/10.1103/PhysRevLett.112.091304}{doi:10.1103/PhysRevLett.112.091304}.

\bibitem{LinYang2004}
F.~Lin and Y.~Yang,
\textit{Existence of energy minimizers as stable knotted solitons in the Faddeev model},
Commun.\ Math.\ Phys.\ \textbf{249}, 273--303 (2004).
\href{https://doi.org/10.1007/s00220-004-1110-y}{doi:10.1007/s00220-004-1110-y}.

\bibitem{Wetterich2024}
C.~Wetterich,
\textit{Dark energy evolution from quantum gravity},
arXiv:2407.03465 [hep-th] (2024).

\bibitem{VasakCCGG2020}
D.~Vasak, J.~Struckmeier, and J.~Kirsch,
\textit{Dark energy and inflation invoked in CCGG by locally contorted space-time},
Eur.\ Phys.\ J.\ Plus \textbf{135}, 404 (2020).
\href{https://doi.org/10.1140/epjp/s13360-020-00415-7}{doi:10.1140/epjp/s13360-020-00415-7}.

\end{thebibliography}

\end{document}
